% Complete reviewed Unit 3; editorial master: chapters/ch03/student.md
% Markdown SHA-256: 0c00a1f1d96fe7a4cbc11b9b633021a9c0cb7a73c1296e44a0d74a7d3da28b49
% One-use initial assembly followed by complete content and page review.
\newcommand{\runningtitle}{Unit 3. Adjectives and number}
% Layout inherited from the reviewed previous course; set a running title per unit.
% This is a small hand-maintained TeX layout, not a conversion pipeline.
\RequirePackage{fix-cm}
\documentclass[11pt,twoside]{article}
\usepackage{fontspec}
\setmainfont{Cambria}[
  Renderer=HarfBuzz,
  BoldFont=cambriab.ttf,
  ItalicFont=cambriai.ttf,
  BoldItalicFont=cambriaz.ttf
]
\setmonofont{lmmono10-regular.otf}[BoldFont=lmmonolt10-bold.otf]
\usepackage[paperwidth=165mm,paperheight=241mm,inner=21mm,outer=19mm,
  top=22mm,bottom=20mm,headheight=14pt,headsep=7mm,footskip=10mm]{geometry}
\usepackage{xcolor,array,booktabs,longtable,tabularx}
\usepackage{enumitem,fancyhdr}
\usepackage[unicode,hidelinks]{hyperref}
\definecolor{ink}{HTML}{000000}
\definecolor{muted}{HTML}{4C5558}
\setlength{\parindent}{0pt}
\setlength{\parskip}{3.5pt}
\AtBeginDocument{\fontsize{11}{14.3}\selectfont}
\setlength{\tabcolsep}{4pt}
\renewcommand{\arraystretch}{1.10}
\setlist[itemize]{leftmargin=1.3em,label=\textbullet,itemsep=2pt,topsep=3pt,parsep=0pt}
\setlist[enumerate]{leftmargin=1.8em,itemsep=3pt,topsep=4pt,parsep=0pt}
\widowpenalty=10000
\clubpenalty=10000
\emergencystretch=1em
\raggedbottom
\pagestyle{fancy}
\fancyhf{}
\providecommand{\runningtitle}{Old English}
\fancyhead[LE,RO]{\small\itshape \runningtitle}
\fancyhead[LO,RE]{\small Old English}
\fancyfoot[LE,RO]{\small\thepage}
\renewcommand{\headrulewidth}{0.3pt}
\newcommand{\booktitle}[2]{
  {\fontsize{9}{11}\selectfont\color{muted}#1\par}
  \vspace{3pt}
  {\fontsize{20}{24}\selectfont\color{ink}#2\par}
  \vspace{6pt}
}
\newcommand{\source}[1]{\par{\footnotesize\color{muted}#1\par}}
\newcommand{\readingtext}[1]{
  \par\addvspace{5pt}
  \begingroup\leftskip=1.5em
    \fontsize{12}{16}\selectfont #1\par
  \endgroup\addvspace{5pt}
}
\makeatletter
\renewcommand{\section}{\@startsection{section}{1}{\z@}
  {-2.4ex \@plus -0.6ex \@minus -0.2ex}{0.8ex \@plus 0.2ex}
  {\normalfont\fontsize{12}{15}\selectfont\bfseries\color{ink}}}
\renewcommand{\subsection}{\@startsection{subsection}{2}{\z@}
  {-1.8ex \@plus -0.5ex \@minus -0.2ex}{0.5ex \@plus 0.2ex}
  {\normalfont\fontsize{11}{14}\selectfont\bfseries\color{ink}}}
\makeatother

\hyphenpenalty=10000
\exhyphenpenalty=10000
\emergencystretch=3em
\newcommand{\itempara}[2]{\par\addvspace{3pt}\noindent\hangindent=1.8em\hangafter=1\makebox[1.8em][l]{#1.}#2\par}
\newcommand{\reserve}[1]{\par\penalty-100\begingroup\dimen0=\pagegoal\advance\dimen0 by-\pagetotal\ifdim\dimen0<#1\newpage\fi\endgroup}
\begin{document}
\raggedright
\booktitle{UNIT 3}{Adjectives, comparison and number}

% mdparagraph:P001:begin
This unit adds description, possession, comparison and counting to the noun and indicative-verb systems of Units 1–2. The grammar and vocabulary below supply the new material needed for independent practice. The final source passage has additional forms and constructions; its local help remains available throughout the reading tasks.
% mdparagraph:P001:end

% mdparagraph:P002:begin
In tables, \textbf{m.}, \textbf{n.} and \textbf{f.} mean masculine, neuter and feminine; \textbf{sg.} and \textbf{pl.} mean singular and plural. \textbf{N/\allowbreak{}G/\allowbreak{}D/\allowbreak{}A} mean nominative, genitive, dative and accusative. Use the selected forms when a task requests them. A documented alternative may be valid Old English even when it does not meet that particular instruction. Length marks and the distinction between e and ę are retained.
% mdparagraph:P002:end

\reserve{7\baselineskip}
% mdparagraph:P003:begin
\section*{1. Descriptive vocabulary}
% mdparagraph:P003:end

% teaching: NEW-U03-VOC-god-adjective
% teaching: NEW-U03-VOC-glaed
% teaching: NEW-U03-VOC-micel
% teaching: NEW-U03-VOC-idel
% teaching: NEW-U03-VOC-aemtig
% teaching: NEW-U03-VOC-aeppelbaere
% teaching: NEW-U03-VOC-lytel
% teaching: NEW-U03-VOC-styriendlic
% teaching: NEW-U03-VOC-leof
% teaching: NEW-U03-VOC-eald
% teaching: NEW-U03-VOC-yfel
% teaching: NEW-U03-VOC-gearu
% teaching: NEW-U03-VOC-heah
% mdtable:T01:begin
\par\begingroup
\fontsize{10.5}{13.2}\selectfont
\setlength{\LTleft}{0pt}\setlength{\LTright}{\fill}
\noindent
\begin{tabular}{@{}>{\raggedright\arraybackslash}p{31.04mm}>{\raggedright\arraybackslash}p{36.68mm}>{\raggedright\arraybackslash}p{50.65mm}@{}}
\toprule
\textbf{Adjective} & \textbf{Meaning} & \textbf{Inflection guide} \\
\midrule
gōd & good & long-stem model \\
glæd & glad & short stem; æ/\allowbreak{}a \\
micel & great; large & micl-/\allowbreak{}micel- \\
īdel & empty; void & īdl-/\allowbreak{}īdel- \\
ǣmtig & empty; void & retain ǣmtig- \\
æppelbǣre & fruit-bearing & -e adjective; æppelbǣr- \\
lȳtel & little; small & lȳtl-/\allowbreak{}lȳtel- \\
styriendlic & moving & regular long stem \\
lēof & dear & regular long stem \\
eald & old & regular long stem \\
yfel & bad; evil & selected full yfel- stem \\
gearu & ready & gearw-/\allowbreak{}gearo- \\
hēah & high & special contraction \\
\bottomrule
\end{tabular}
\par\endgroup

% mdtable:T01:end
% mdparagraph:P004:begin
Learn the meaning with the inflection guide. \textbf{Styriendlic} is an adjective meaning “moving”; its use here does not require a participle paradigm. The particular compound \textbf{æppelbǣre} is Cook’s word for “fruit-bearing”. The following tables and stem map give the complete selected forms of every adjective in this list.
% mdparagraph:P004:end

\reserve{7\baselineskip}
% mdparagraph:P005:begin
\section*{2. Agreement and the two declensions}
% mdparagraph:P005:end

% teaching: NEW-U03-GRAM-agreement
% teaching: NEW-U03-GRAM-determiner-conditions
% teaching: NEW-U03-GRAM-predicate-adjective
% mdparagraph:P006:begin
An adjective agrees with its own noun in gender, number and case. An attributive adjective belongs inside a noun phrase. A predicate adjective states something about the subject with the copula and is nominative. Agreement is grammatical: a neuter noun takes a neuter adjective regardless of the English meaning.
% mdparagraph:P006:end

% mdparagraph:P007:begin
Ordinary adjectives have a \textbf{strong} and a \textbf{weak} declension. Use the strong form in a bare noun phrase and for an ordinary positive predicate adjective. Use the weak form after the article or a demonstrative. An article inside the subject does not make its predicate adjective weak. Comparatives are always weak; ordinals and superlatives have the rules given later.
% mdparagraph:P007:end

% mdparagraph:P008:begin
The following examples are newly composed teaching material.
% mdparagraph:P008:end

% example: NEW-U03-EXAMPLE-01
% mdparagraph:P009:begin
Gōd męte is on ðǣre eorðan.
% mdparagraph:P009:end

% mdparagraph:P010:begin
Good food is on the earth.
% mdparagraph:P010:end

% mdparagraph:P011:begin
Gōd agrees with the masculine nominative singular męte. The locative phrase retains the dative taught in Unit 1.
% mdparagraph:P011:end

% example: NEW-U03-EXAMPLE-02
% mdparagraph:P012:begin
Se gōda mann hæfð gōdne męte.
% mdparagraph:P012:end

% mdparagraph:P013:begin
The good man has good food.
% mdparagraph:P013:end

% mdparagraph:P014:begin
Gōda is weak after se and agrees with the subject mann. Gōdne is strong masculine accusative singular and agrees with the bare object męte.
% mdparagraph:P014:end

% example: NEW-U03-EXAMPLE-03
% mdparagraph:P015:begin
Ðæt gōde trēo is micel.
% mdparagraph:P015:end

% mdparagraph:P016:begin
The good tree is great.
% mdparagraph:P016:end

% mdparagraph:P017:begin
Gōde is weak inside the determined subject phrase. Micel is a strong neuter nominative singular predicate.
% mdparagraph:P017:end

% mdparagraph:P018:begin
A following ordinary adjective may be strong or weak after a possessive: Cook describes the weak form as occurring sometimes. Do not turn possession into an invariable weak-adjective rule. The exercises avoid requiring a new general rule for an adjective after a limiting noun genitive.
% mdparagraph:P018:end

\clearpage
% mdparagraph:P019:begin
\section*{3. Complete strong and weak models}
% mdparagraph:P019:end

% teaching: NEW-U03-GRAM-strong-adjectives
% teaching: NEW-U03-GRAM-long-stems
\reserve{7\baselineskip}
% mdparagraph:P020:begin
\subsection*{Strong gōd}
% mdparagraph:P020:end

% teaching: NEW-U03-PAR-god-adjective-strong
% mdtable:T02:begin
% course-table-style: 2026-10
% course-panel:T02-P01:begin
\par\begingroup
\fontsize{10.5}{13.2}\selectfont
\setlength{\LTleft}{0pt}\setlength{\LTright}{\fill}
\noindent
\begin{tabular}{@{}>{\raggedright\arraybackslash}p{22.0mm}>{\raggedright\arraybackslash}p{31.188mm}>{\raggedright\arraybackslash}p{31.188mm}>{\raggedright\arraybackslash}p{31.188mm}@{}}
\toprule
\textbf{Case} & \textbf{Masculine singular} & \textbf{Neuter singular} & \textbf{Feminine singular} \\
\midrule
Nominative & gōd & gōd & gōd \\
Genitive & gōdes & gōdes & gōdre \\
Dative & gōdum & gōdum & gōdre \\
Accusative & gōdne & gōd & gōde \\
\bottomrule
\end{tabular}
\par\endgroup
% course-panel:T02-P01:end
\par\addvspace{5pt}
% course-panel:T02-P02:begin
\par\begingroup
\fontsize{10.5}{13.2}\selectfont
\setlength{\LTleft}{0pt}\setlength{\LTright}{\fill}
\noindent
\begin{tabular}{@{}>{\raggedright\arraybackslash}p{22.0mm}>{\raggedright\arraybackslash}p{31.188mm}>{\raggedright\arraybackslash}p{31.188mm}>{\raggedright\arraybackslash}p{31.188mm}@{}}
\toprule
\textbf{Case} & \textbf{Masculine plural} & \textbf{Neuter plural} & \textbf{Feminine plural} \\
\midrule
Nominative & gōde & gōd & gōda \\
Genitive & gōdra & gōdra & gōdra \\
Dative & gōdum & gōdum & gōdum \\
Accusative & gōde & gōd & gōda \\
\bottomrule
\end{tabular}
\par\endgroup
% course-panel:T02-P02:end
% mdtable:T02:end
% mdparagraph:P021:begin
Instrumental singular, recognition only: masculine gōde; neuter gōde.
% mdparagraph:P021:end

% mdparagraph:P022:begin
The endingless feminine nominative singular and neuter nominative/\allowbreak{}accusative plural belong to the selected long-stem pattern. Masculine and feminine plurals need their own endings even when the English adjective is unchanged. Neuter nominative and accusative coincide in both numbers; they remain different cases.
% mdparagraph:P022:end

% teaching: NEW-U03-GRAM-weak-adjectives
\reserve{7\baselineskip}
% mdparagraph:P023:begin
\subsection*{Weak gōd}
% mdparagraph:P023:end

% teaching: NEW-U03-PAR-god-adjective-weak
% mdtable:T03:begin
% course-table-style: 2026-10
% course-panel:T03-P01:begin
\par\begingroup
\fontsize{10.5}{13.2}\selectfont
\setlength{\LTleft}{0pt}\setlength{\LTright}{\fill}
\noindent
\begin{tabular}{@{}>{\raggedright\arraybackslash}p{22.0mm}>{\raggedright\arraybackslash}p{31.188mm}>{\raggedright\arraybackslash}p{31.188mm}>{\raggedright\arraybackslash}p{31.188mm}@{}}
\toprule
\textbf{Case} & \textbf{Masculine singular} & \textbf{Neuter singular} & \textbf{Feminine singular} \\
\midrule
Nominative & gōda & gōde & gōde \\
Genitive & gōdan & gōdan & gōdan \\
Dative & gōdan & gōdan & gōdan \\
Accusative & gōdan & gōde & gōdan \\
\bottomrule
\end{tabular}
\par\endgroup
% course-panel:T03-P01:end
\par\addvspace{5pt}
% course-panel:T03-P02:begin
\par\begingroup
\fontsize{10.5}{13.2}\selectfont
\setlength{\LTleft}{0pt}\setlength{\LTright}{\fill}
\noindent
\begin{tabular}{@{}>{\raggedright\arraybackslash}p{22.0mm}>{\raggedright\arraybackslash}p{31.188mm}>{\raggedright\arraybackslash}p{31.188mm}>{\raggedright\arraybackslash}p{31.188mm}@{}}
\toprule
\textbf{Case} & \textbf{Masculine plural} & \textbf{Neuter plural} & \textbf{Feminine plural} \\
\midrule
Nominative & gōdan & gōdan & gōdan \\
Genitive & gōdra & gōdra & gōdra \\
Dative & gōdum & gōdum & gōdum \\
Accusative & gōdan & gōdan & gōdan \\
\bottomrule
\end{tabular}
\par\endgroup
% course-panel:T03-P02:end
% mdtable:T03:end
% mdparagraph:P024:begin
The weak adjective has -a in masculine nominative singular, -e in feminine nominative singular and neuter nominative/\allowbreak{}accusative singular, and -an in the other singular cells. All genders share the plural forms. Its selected genitive plural is -ra; do not replace it mechanically with a weak noun’s -ena. Cook also admits weak-adjective genitive plurals in -a and -ena. The comparative exceptions below are learned explicitly.
% mdparagraph:P024:end

% teaching: NEW-U03-GRAM-instrumental-reference
% mdparagraph:P025:begin
The distinct strong instrumental singular is recorded for recognition. It has the same spelling as several ordinary -e forms, so a spelling alone does not identify its case. No instrumental phrase or construction is required in independent production here. Weak instrumental forms coincide with the corresponding datives; no separate instrumental use is being introduced.
% mdparagraph:P025:end

\reserve{7\baselineskip}
% mdparagraph:P026:begin
\section*{4. Stems and important exceptions}
% mdparagraph:P026:end

% teaching: NEW-U03-GRAM-short-stems
% teaching: NEW-U03-GRAM-ae-a
\reserve{7\baselineskip}
% mdparagraph:P027:begin
\subsection*{Short-stem glæd}
% mdparagraph:P027:end

% teaching: NEW-U03-PAR-glaed-strong
% mdtable:T04:begin
% course-table-style: 2026-10
% course-panel:T04-P01:begin
\par\begingroup
\fontsize{10.5}{13.2}\selectfont
\setlength{\LTleft}{0pt}\setlength{\LTright}{\fill}
\noindent
\begin{tabular}{@{}>{\raggedright\arraybackslash}p{22.0mm}>{\raggedright\arraybackslash}p{31.188mm}>{\raggedright\arraybackslash}p{31.188mm}>{\raggedright\arraybackslash}p{31.188mm}@{}}
\toprule
\textbf{Case} & \textbf{Masculine singular} & \textbf{Neuter singular} & \textbf{Feminine singular} \\
\midrule
Nominative & glæd & glæd & gladu \\
Genitive & glades & glades & glædre \\
Dative & gladum & gladum & glædre \\
Accusative & glædne & glæd & glade \\
\bottomrule
\end{tabular}
\par\endgroup
% course-panel:T04-P01:end
\par\addvspace{5pt}
% course-panel:T04-P02:begin
\par\begingroup
\fontsize{10.5}{13.2}\selectfont
\setlength{\LTleft}{0pt}\setlength{\LTright}{\fill}
\noindent
\begin{tabular}{@{}>{\raggedright\arraybackslash}p{22.0mm}>{\raggedright\arraybackslash}p{31.188mm}>{\raggedright\arraybackslash}p{31.188mm}>{\raggedright\arraybackslash}p{31.188mm}@{}}
\toprule
\textbf{Case} & \textbf{Masculine plural} & \textbf{Neuter plural} & \textbf{Feminine plural} \\
\midrule
Nominative & glade & gladu & glada \\
Genitive & glædra & glædra & glædra \\
Dative & gladum & gladum & gladum \\
Accusative & glade & gladu & glada \\
\bottomrule
\end{tabular}
\par\endgroup
% course-panel:T04-P02:end
% mdtable:T04:end
% mdparagraph:P028:begin
Instrumental singular, recognition only: masculine glade; neuter glade.
% mdparagraph:P028:end

% teaching: NEW-U03-PAR-glaed-weak
% mdparagraph:P029:begin
The selected short-stem feminine nominative singular and neuter nominative/\allowbreak{}accusative plural have -u. Glæd uses glad- before the vowel endings and glæd- before -ne, -re and -ra. Its complete weak declension is the weak gōd pattern with glad- before -a, -e, -an and -um, but glæd- before genitive plural -ra. This gives a complete pattern; it is not a rule for changing æ to a in every adjective.
% mdparagraph:P029:end

% teaching: NEW-U03-GRAM-syncope
% teaching: NEW-U03-GRAM-e-stem
% teaching: NEW-U03-GRAM-u-w
\reserve{7\baselineskip}
% mdparagraph:P030:begin
\subsection*{A complete application map}
% mdparagraph:P030:end

% mdparagraph:P031:begin
Use the strong and weak gōd panels with this map. \textbf{V-stem} is used before endings beginning with a vowel: -es, -um, -e, -a and -an. \textbf{C-stem} is used before -ne, -re and -ra. The strong masculine nominative singular and neuter nominative/\allowbreak{}accusative singular are the listed headword. The last two columns replace the corresponding endingless gōd cells. The strong instrumental singular is V-stem + -e. Thus every strong and weak N/\allowbreak{}G/\allowbreak{}D/\allowbreak{}A cell is determined; do not add an ending twice to a form already listed in the last columns.
% mdparagraph:P031:end

% teaching: NEW-U03-PAR-micel-strong
% teaching: NEW-U03-PAR-micel-weak
% teaching: NEW-U03-PAR-idel-strong
% teaching: NEW-U03-PAR-idel-weak
% teaching: NEW-U03-PAR-aemtig-strong
% teaching: NEW-U03-PAR-aemtig-weak
% teaching: NEW-U03-PAR-aeppelbaere-strong
% teaching: NEW-U03-PAR-aeppelbaere-weak
% teaching: NEW-U03-PAR-lytel-strong
% teaching: NEW-U03-PAR-lytel-weak
% teaching: NEW-U03-PAR-styriendlic-strong
% teaching: NEW-U03-PAR-styriendlic-weak
% teaching: NEW-U03-PAR-leof-strong
% teaching: NEW-U03-PAR-leof-weak
% teaching: NEW-U03-PAR-eald-strong
% teaching: NEW-U03-PAR-eald-weak
% teaching: NEW-U03-PAR-yfel-strong
% teaching: NEW-U03-PAR-yfel-weak
% teaching: NEW-U03-PAR-gearu-strong
% teaching: NEW-U03-PAR-gearu-weak
% mdtable:T05:begin
\par\begingroup
\fontsize{10.5}{13.2}\selectfont
\setlength{\LTleft}{0pt}\setlength{\LTright}{\fill}
\noindent
\begin{tabular}{@{}>{\raggedright\arraybackslash}p{19.73mm}>{\raggedright\arraybackslash}p{21.53mm}>{\raggedright\arraybackslash}p{21.53mm}>{\raggedright\arraybackslash}p{25.38mm}>{\raggedright\arraybackslash}p{24.58mm}@{}}
\toprule
\textbf{Headword} & \textbf{V-stem} & \textbf{C-stem} & \textbf{Strong f. nom. sg.} & \textbf{Strong n. N/\allowbreak{}A pl.} \\
\midrule
gōd & gōd- & gōd- & gōd & gōd \\
glæd & glad- & glæd- & gladu & gladu \\
micel & micl- & micel- & miclu & miclu \\
īdel & īdl- & īdel- & īdel & īdlu \\
ǣmtig & ǣmtig- & ǣmtig- & ǣmtig & ǣmtig \\
æppelbǣre & æppelbǣr- & æppelbǣr- & æppelbǣru & æppelbǣru \\
lȳtel & lȳtl- & lȳtel- & lȳtlu & lȳtlu \\
styriendlic & styriendlic- & styriendlic- & styriendlic & styriendlic \\
lēof & lēof- & lēof- & lēof & lēof \\
eald & eald- & eald- & eald & eald \\
yfel & yfel- & yfel- & yfel & yfel \\
gearu & gearw- & gearo- & gearu & gearu \\
\bottomrule
\end{tabular}
\par\endgroup

% mdtable:T05:end
% mdparagraph:P032:begin
In micel, lȳtel and īdel, the medial vowel is lost before the selected vowel endings but retained before -ne, -re and -ra. The selected īdel feminine nominative singular is endingless, while its neuter plural is īdlu. Retain the full stem in ǣmtig and in the selected yfel pattern. A long disyllable does not license automatic deletion of every medial vowel.
% mdparagraph:P032:end

% mdparagraph:P033:begin
In æppelbǣre the final headword -e is replaced by the endings, rather than treated as part of the stem. Its C-stem ends in r: adding -re or -ra therefore writes rr. Gearu has w before vowel endings but the selected gearo- before consonantal endings; its listed direct forms remain gearu. The map includes the weak declensions of both words.
% mdparagraph:P033:end

% mdparagraph:P034:begin
Selected alternatives worth recognizing: strong feminine nominative micel or micelu beside miclu; lȳtel or lȳtelu beside lȳtlu; īdelu beside īdel. Corresponding neuter plurals may retain or omit -u, and strong feminine N/\allowbreak{}A plurals often have -e beside selected -a. These alternatives do not permit changing the actual spelling of a source passage. In particular, the reading prints miclu.
% mdparagraph:P034:end

% teaching: NEW-U03-GRAM-h-loss
\reserve{7\baselineskip}
% mdparagraph:P035:begin
\subsection*{Hēah}
% mdparagraph:P035:end

% mdparagraph:P036:begin
Hēah loses h and contracts in much of its inflection. Use this complete selected pattern instead of attaching the ordinary endings directly to the headword.
% mdparagraph:P036:end

% teaching: NEW-U03-PAR-heah-strong
\par\begin{minipage}{\linewidth}\raggedright
% mdparagraph:P037:begin
\textbf{Strong hēah.}
% mdparagraph:P037:end

% mdtable:T06:begin
% course-table-style: 2026-10
% course-panel:T06-P01:begin
\par\begingroup
\fontsize{10.5}{13.2}\selectfont
\setlength{\LTleft}{0pt}\setlength{\LTright}{\fill}
\noindent
\begin{tabular}{@{}>{\raggedright\arraybackslash}p{22.0mm}>{\raggedright\arraybackslash}p{31.188mm}>{\raggedright\arraybackslash}p{31.188mm}>{\raggedright\arraybackslash}p{31.188mm}@{}}
\toprule
\textbf{Case} & \textbf{Masculine singular} & \textbf{Neuter singular} & \textbf{Feminine singular} \\
\midrule
Nominative & hēah & hēah & hēa \\
Genitive & hēas & hēas & hēarre \\
Dative & hēaum & hēaum & hēarre \\
Accusative & hēanne & hēah & hēa \\
\bottomrule
\end{tabular}
\par\endgroup
% course-panel:T06-P01:end
\par\addvspace{5pt}
% course-panel:T06-P02:begin
\par\begingroup
\fontsize{10.5}{13.2}\selectfont
\setlength{\LTleft}{0pt}\setlength{\LTright}{\fill}
\noindent
\begin{tabular}{@{}>{\raggedright\arraybackslash}p{22.0mm}>{\raggedright\arraybackslash}p{31.188mm}>{\raggedright\arraybackslash}p{31.188mm}>{\raggedright\arraybackslash}p{31.188mm}@{}}
\toprule
\textbf{Case} & \textbf{Masculine plural} & \textbf{Neuter plural} & \textbf{Feminine plural} \\
\midrule
Nominative & hēa & hēa & hēa \\
Genitive & hēarra & hēarra & hēarra \\
Dative & hēaum & hēaum & hēaum \\
Accusative & hēa & hēa & hēa \\
\bottomrule
\end{tabular}
\par\endgroup
% course-panel:T06-P02:end
% mdtable:T06:end
\end{minipage}\par
% mdparagraph:P038:begin
Instrumental singular, recognition only: masculine hēa; neuter hēa.
% mdparagraph:P038:end

% teaching: NEW-U03-PAR-heah-weak
\par\begin{minipage}{\linewidth}\raggedright
% mdparagraph:P039:begin
\textbf{Weak hēah.}
% mdparagraph:P039:end

% mdtable:T07:begin
% course-table-style: 2026-10
% course-panel:T07-P01:begin
\par\begingroup
\fontsize{10.5}{13.2}\selectfont
\setlength{\LTleft}{0pt}\setlength{\LTright}{\fill}
\noindent
\begin{tabular}{@{}>{\raggedright\arraybackslash}p{22.0mm}>{\raggedright\arraybackslash}p{31.188mm}>{\raggedright\arraybackslash}p{31.188mm}>{\raggedright\arraybackslash}p{31.188mm}@{}}
\toprule
\textbf{Case} & \textbf{Masculine singular} & \textbf{Neuter singular} & \textbf{Feminine singular} \\
\midrule
Nominative & hēa & hēa & hēa \\
Genitive & hēan & hēan & hēan \\
Dative & hēan & hēan & hēan \\
Accusative & hēan & hēa & hēan \\
\bottomrule
\end{tabular}
\par\endgroup
% course-panel:T07-P01:end
\par\addvspace{5pt}
% course-panel:T07-P02:begin
\par\begingroup
\fontsize{10.5}{13.2}\selectfont
\setlength{\LTleft}{0pt}\setlength{\LTright}{\fill}
\noindent
\begin{tabular}{@{}>{\raggedright\arraybackslash}p{22.0mm}>{\raggedright\arraybackslash}p{31.188mm}>{\raggedright\arraybackslash}p{31.188mm}>{\raggedright\arraybackslash}p{31.188mm}@{}}
\toprule
\textbf{Case} & \textbf{Masculine plural} & \textbf{Neuter plural} & \textbf{Feminine plural} \\
\midrule
Nominative & hēan & hēan & hēan \\
Genitive & hēarra & hēarra & hēarra \\
Dative & hēaum & hēaum & hēaum \\
Accusative & hēan & hēan & hēan \\
\bottomrule
\end{tabular}
\par\endgroup
% course-panel:T07-P02:end
% mdtable:T07:end
\end{minipage}\par
% mdparagraph:P040:begin
The selected hēaum has the supported alternative hēam. Hēanne may appear as hēane; hēarre/\allowbreak{}hēarra as hēare/\allowbreak{}hēara. The tables retain their selected doubling and vowel sequence. These forms show why consonant loss, doubling and contraction must be learned with the paradigm.
% mdparagraph:P040:end

% teaching: NEW-U03-VOC-eall
% teaching: NEW-U03-PAR-eall-strong
\reserve{7\baselineskip}
% mdparagraph:P041:begin
\subsection*{Reference: eall}
% mdparagraph:P041:end

% mdparagraph:P042:begin
Eall means “all”. Its strong morphology is for recognition here; no independent quantifier construction is required. Apply the complete strong gōd pattern with headword eall, V-stem eall-, C-stem eal-, feminine nominative singular eall and neuter N/\allowbreak{}A plural eall. The instrumental singular is ealle. The single l before -ne/\allowbreak{}-re/\allowbreak{}-ra is deliberate. Do not construct a weak paradigm from this reference.
% mdparagraph:P042:end

\reserve{7\baselineskip}
% mdparagraph:P043:begin
\section*{5. Things and counted lights}
% mdparagraph:P043:end

% teaching: NEW-U03-VOC-thing
% teaching: NEW-U03-VOC-leoht
% teaching: NEW-U03-GRAM-thing-noun
% teaching: NEW-U03-GRAM-light-count-sense
% teaching: NEW-U03-PAR-thing
% teaching: NEW-U03-PAR-leoht-count
% mdparagraph:P044:begin
Ðing, neuter, means “thing”. Lēoht, neuter, can now also name an individual light or luminary. This counted sense differs from the mass sense “light” already learned. Both follow the complete heavy-neuter pattern:
% mdparagraph:P044:end

% mdtable:T08:begin
% course-table-style: 2026-10
% course-panel:T08-P01:begin
\par\begingroup
\fontsize{10.5}{13.2}\selectfont
\setlength{\LTleft}{0pt}\setlength{\LTright}{\fill}
\noindent
\begin{tabular}{@{}>{\raggedright\arraybackslash}p{22.0mm}>{\raggedright\arraybackslash}p{22.688mm}>{\raggedright\arraybackslash}p{22.688mm}>{\raggedright\arraybackslash}p{22.688mm}>{\raggedright\arraybackslash}p{22.688mm}@{}}
\toprule
\textbf{Case} & \textbf{ðing singular} & \textbf{ðing plural} & \textbf{lēoht singular} & \textbf{lēoht plural} \\
\midrule
Nominative & ðing & ðing & lēoht & lēoht \\
Genitive & ðinges & ðinga & lēohtes & lēohta \\
Dative & ðinge & ðingum & lēohte & lēohtum \\
Accusative & ðing & ðing & lēoht & lēoht \\
\bottomrule
\end{tabular}
\par\endgroup
% course-panel:T08-P01:end
% mdtable:T08:end
% mdparagraph:P045:begin
Use all four plural cases of counted lēoht productively. The unchanged spelling lēoht in nominative/\allowbreak{}accusative plural does not make it singular. Number is also shown by its determiner, numeral, adjective and clause context.
% mdparagraph:P045:end

\reserve{7\baselineskip}
% mdparagraph:P046:begin
\subsection*{Adjectives — A1–A10}
% mdparagraph:P046:end

\par\begin{minipage}{\linewidth}\raggedright
% mdparagraph:P047:begin
\textbf{A1.}
% mdparagraph:P047:end

% item: NEW-U03-A01
% mdparagraph:P048:begin
Give the complete strong nominative, genitive, dative and accusative of gōd in all three genders, singular and plural. Label every case even where spellings coincide.
% mdparagraph:P048:end

\end{minipage}\par\addvspace{4pt}

\par\begin{minipage}{\linewidth}\raggedright
% mdparagraph:P049:begin
\textbf{A2.}
% mdparagraph:P049:end

% item: NEW-U03-A02
% mdparagraph:P050:begin
Give the complete strong nominative, genitive, dative and accusative of glæd in all three genders, singular and plural. Mark the forms with a rather than æ.
% mdparagraph:P050:end

\end{minipage}\par\addvspace{4pt}

\par\begin{minipage}{\linewidth}\raggedright
% mdparagraph:P051:begin
\textbf{A3.}
% mdparagraph:P051:end

% item: NEW-U03-A03
% mdparagraph:P052:begin
Give the complete weak nominative, genitive, dative and accusative of gōd in all three genders, singular and plural. Contrast its genitive plural with a weak noun’s genitive plural.
% mdparagraph:P052:end

\end{minipage}\par\addvspace{4pt}

\par\begin{minipage}{\linewidth}\raggedright
% mdparagraph:P053:begin
\textbf{A4.}
% mdparagraph:P053:end

% item: NEW-U03-A04
% mdparagraph:P054:begin
For micel, lȳtel, īdel and ǣmtig, give the strong masculine genitive singular, strong feminine nominative singular, strong neuter nominative plural and weak neuter nominative singular. Preserve the selected stems and vowel marks.
% mdparagraph:P054:end

\end{minipage}\par\addvspace{4pt}

\par\begin{minipage}{\linewidth}\raggedright
% mdparagraph:P055:begin
\textbf{A5.}
% mdparagraph:P055:end

% item: NEW-U03-A05
% mdparagraph:P056:begin
For æppelbǣre, give the strong masculine and feminine nominative singular, masculine accusative singular, neuter nominative plural and genitive plural. Explain why the headword ends in -e. For gearu and hēah, give the strong masculine nominative, genitive, dative and accusative singular and the weak masculine nominative singular.
% mdparagraph:P056:end

\end{minipage}\par\addvspace{4pt}

\par\begin{minipage}{\linewidth}\raggedright
% mdparagraph:P057:begin
\textbf{A6.}
% mdparagraph:P057:end

% item: NEW-U03-A06
% mdparagraph:P058:begin
Give the four cases of ðing in the singular and plural. Give its meaning in the plural phrase ðā ðing.
% mdparagraph:P058:end

\end{minipage}\par\addvspace{4pt}

\par\begin{minipage}{\linewidth}\raggedright
% mdparagraph:P059:begin
\textbf{A7.}
% mdparagraph:P059:end

% item: NEW-U03-A07
% mdparagraph:P060:begin
These phrases are nominative. Add the appropriate article to each and change the adjective as needed: gōd fugol; gōd wyrt; gōd trēo; gōde fuglas. Keep the same case and number.
% mdparagraph:P060:end

\end{minipage}\par\addvspace{4pt}

\par\begin{minipage}{\linewidth}\raggedright
% mdparagraph:P061:begin
\textbf{A8.}
% mdparagraph:P061:end

% item: NEW-U03-A08
% mdparagraph:P062:begin
Translate and explain the different adjective forms: Ðā gōdan fuglas sindon on ðǣm hofe. Ðā fuglas sindon gōde.
% mdparagraph:P062:end

\end{minipage}\par\addvspace{4pt}

\par\begin{minipage}{\linewidth}\raggedright
% mdparagraph:P063:begin
\textbf{A9.}
% mdparagraph:P063:end

% item: NEW-U03-A09
% mdparagraph:P064:begin
Translate “The birds live in the good man’s dwelling.” Put the limiting genitive before the head noun. Identify the case of the dwelling and of the man.
% mdparagraph:P064:end

\end{minipage}\par\addvspace{4pt}

\par\begin{minipage}{\linewidth}\raggedright
% mdparagraph:P065:begin
\textbf{A10.}
% mdparagraph:P065:end

% item: NEW-U03-A10
% mdparagraph:P066:begin
The strong table contains gōde as masculine nominative plural and as masculine or neuter instrumental singular. State each analysis. Compare the corresponding entries of glæd.
% mdparagraph:P066:end

\end{minipage}\par\addvspace{4pt}

\reserve{7\baselineskip}
% mdparagraph:P067:begin
\section*{6. Agreeing possessives}
% mdparagraph:P067:end

% teaching: NEW-U03-GRAM-possessive-agreement
% teaching: NEW-U03-GRAM-possessor-reference
% teaching: NEW-U03-VOC-min-poss
% teaching: NEW-U03-VOC-thin-poss
% teaching: NEW-U03-VOC-uncer-poss
% teaching: NEW-U03-VOC-incer-poss
% teaching: NEW-U03-VOC-ure-poss
% teaching: NEW-U03-VOC-eower-poss
% mdtable:T09:begin
\par\begingroup
\fontsize{10.5}{13.2}\selectfont
\setlength{\LTleft}{0pt}\setlength{\LTright}{\fill}
\noindent
\begin{tabular}{@{}>{\raggedright\arraybackslash}p{47.38mm}>{\raggedright\arraybackslash}p{39.74mm}>{\raggedright\arraybackslash}p{31.26mm}@{}}
\toprule
\textbf{Possessor} & \textbf{Agreeing possessive} & \textbf{Meaning} \\
\midrule
the speaker & mīn & my \\
one addressee & ðīn & your \\
the speaker and one other & uncer & our, of us two \\
two addressees & incer & your, of you two \\
the speaker and others & ūre & our \\
several addressees & ēower & your \\
\bottomrule
\end{tabular}
\par\endgroup

% mdtable:T09:end
% mdparagraph:P068:begin
The stem identifies the possessor; the ending agrees with the possessed noun. A dual possessor can own one thing or several things. A plural possessive does not require its head noun to be plural. These agreeing adjectives are distinct from the personal-genitive citation forms recognized in Unit 2. The third-person possessors his, hiere and hiera remain invariant.
% mdparagraph:P068:end

% example: NEW-U03-EXAMPLE-04
% mdparagraph:P069:begin
Ðū hæfst mīnne fisc.
% mdparagraph:P069:end

% mdparagraph:P070:begin
You have my fish.
% mdparagraph:P070:end

% mdparagraph:P071:begin
The possessor is the speaker. The possessed object is masculine accusative singular, so the selected possessive is mīnne.
% mdparagraph:P071:end

% example: NEW-U03-EXAMPLE-05
% mdparagraph:P072:begin
Hīe hīerað ūre fuglas.
% mdparagraph:P072:end

% mdparagraph:P073:begin
They hear our birds.
% mdparagraph:P073:end

% mdparagraph:P074:begin
The possessors are the speaker and others; the birds are masculine accusative plural. The possessive therefore has the selected form ūre.
% mdparagraph:P074:end

\reserve{7\baselineskip}
% mdparagraph:P075:begin
\subsection*{Mīn and ðīn}
% mdparagraph:P075:end

% teaching: NEW-U03-PAR-min-poss-strong
% mdtable:T10:begin
% course-table-style: 2026-10
% course-panel:T10-P01:begin
\par\begingroup
\fontsize{10.5}{13.2}\selectfont
\setlength{\LTleft}{0pt}\setlength{\LTright}{\fill}
\noindent
\begin{tabular}{@{}>{\raggedright\arraybackslash}p{22.0mm}>{\raggedright\arraybackslash}p{31.188mm}>{\raggedright\arraybackslash}p{31.188mm}>{\raggedright\arraybackslash}p{31.188mm}@{}}
\toprule
\textbf{Case} & \textbf{Masculine singular} & \textbf{Neuter singular} & \textbf{Feminine singular} \\
\midrule
Nominative & mīn & mīn & mīn \\
Genitive & mīnes & mīnes & mīnre \\
Dative & mīnum & mīnum & mīnre \\
Accusative & mīnne & mīn & mīne \\
\bottomrule
\end{tabular}
\par\endgroup
% course-panel:T10-P01:end
\par\addvspace{5pt}
% course-panel:T10-P02:begin
\par\begingroup
\fontsize{10.5}{13.2}\selectfont
\setlength{\LTleft}{0pt}\setlength{\LTright}{\fill}
\noindent
\begin{tabular}{@{}>{\raggedright\arraybackslash}p{22.0mm}>{\raggedright\arraybackslash}p{31.188mm}>{\raggedright\arraybackslash}p{31.188mm}>{\raggedright\arraybackslash}p{31.188mm}@{}}
\toprule
\textbf{Case} & \textbf{Masculine plural} & \textbf{Neuter plural} & \textbf{Feminine plural} \\
\midrule
Nominative & mīne & mīn & mīna \\
Genitive & mīnra & mīnra & mīnra \\
Dative & mīnum & mīnum & mīnum \\
Accusative & mīne & mīn & mīna \\
\bottomrule
\end{tabular}
\par\endgroup
% course-panel:T10-P02:end
% mdtable:T10:end
% mdparagraph:P076:begin
Instrumental singular, recognition only: masculine mīne; neuter mīne.
% mdparagraph:P076:end

% teaching: NEW-U03-PAR-thin-poss-strong
% mdparagraph:P077:begin
The complete ðīn paradigm is obtained by replacing mīn- with ðīn- in every cell, preserving every ending and doubled n. Its masculine/\allowbreak{}neuter instrumental singular is ðīne. These two long stems do not take the short-stem -u in the selected feminine nominative singular or neuter N/\allowbreak{}A plural.
% mdparagraph:P077:end

\reserve{7\baselineskip}
% mdparagraph:P078:begin
\subsection*{Uncer and incer}
% mdparagraph:P078:end

% teaching: NEW-U03-PAR-uncer-poss-strong
% mdtable:T11:begin
% course-table-style: 2026-10
% course-panel:T11-P01:begin
\par\begingroup
\fontsize{10.5}{13.2}\selectfont
\setlength{\LTleft}{0pt}\setlength{\LTright}{\fill}
\noindent
\begin{tabular}{@{}>{\raggedright\arraybackslash}p{22.0mm}>{\raggedright\arraybackslash}p{31.188mm}>{\raggedright\arraybackslash}p{31.188mm}>{\raggedright\arraybackslash}p{31.188mm}@{}}
\toprule
\textbf{Case} & \textbf{Masculine singular} & \textbf{Neuter singular} & \textbf{Feminine singular} \\
\midrule
Nominative & uncer & uncer & uncru \\
Genitive & uncres & uncres & uncerre \\
Dative & uncrum & uncrum & uncerre \\
Accusative & uncerne & uncer & uncre \\
\bottomrule
\end{tabular}
\par\endgroup
% course-panel:T11-P01:end
\par\addvspace{5pt}
% course-panel:T11-P02:begin
\par\begingroup
\fontsize{10.5}{13.2}\selectfont
\setlength{\LTleft}{0pt}\setlength{\LTright}{\fill}
\noindent
\begin{tabular}{@{}>{\raggedright\arraybackslash}p{22.0mm}>{\raggedright\arraybackslash}p{31.188mm}>{\raggedright\arraybackslash}p{31.188mm}>{\raggedright\arraybackslash}p{31.188mm}@{}}
\toprule
\textbf{Case} & \textbf{Masculine plural} & \textbf{Neuter plural} & \textbf{Feminine plural} \\
\midrule
Nominative & uncre & uncru & uncra \\
Genitive & uncerra & uncerra & uncerra \\
Dative & uncrum & uncrum & uncrum \\
Accusative & uncre & uncru & uncra \\
\bottomrule
\end{tabular}
\par\endgroup
% course-panel:T11-P02:end
% mdtable:T11:end
% mdparagraph:P079:begin
Instrumental singular, recognition only: masculine uncre; neuter uncre.
% mdparagraph:P079:end

% teaching: NEW-U03-PAR-incer-poss-strong
% mdparagraph:P080:begin
The complete incer paradigm replaces uncer-/\allowbreak{}uncr- with incer-/\allowbreak{}incr- in the corresponding cells. Retain the medial e before -ne, -re and -ra; omit it before the selected vowel endings. Its instrumental singular is incre. Unsyncopated unceru/\allowbreak{}inceru are supported alternatives to uncru/\allowbreak{}incru, and single-r forms such as uncere/\allowbreak{}uncera or incere/\allowbreak{}incera occur beside the selected doubled-r forms.
% mdparagraph:P080:end

\reserve{7\baselineskip}
% mdparagraph:P081:begin
\subsection*{Ūre and ēower}
% mdparagraph:P081:end

% teaching: NEW-U03-PAR-ure-poss-strong
% mdtable:T12:begin
% course-table-style: 2026-10
% course-panel:T12-P01:begin
\par\begingroup
\fontsize{10.5}{13.2}\selectfont
\setlength{\LTleft}{0pt}\setlength{\LTright}{\fill}
\noindent
\begin{tabular}{@{}>{\raggedright\arraybackslash}p{22.0mm}>{\raggedright\arraybackslash}p{31.188mm}>{\raggedright\arraybackslash}p{31.188mm}>{\raggedright\arraybackslash}p{31.188mm}@{}}
\toprule
\textbf{Case} & \textbf{Masculine singular} & \textbf{Neuter singular} & \textbf{Feminine singular} \\
\midrule
Nominative & ūre & ūre & ūre \\
Genitive & ūres & ūres & ūrre \\
Dative & ūrum & ūrum & ūrre \\
Accusative & ūrne & ūre & ūre \\
\bottomrule
\end{tabular}
\par\endgroup
% course-panel:T12-P01:end
\par\addvspace{5pt}
% course-panel:T12-P02:begin
\par\begingroup
\fontsize{10.5}{13.2}\selectfont
\setlength{\LTleft}{0pt}\setlength{\LTright}{\fill}
\noindent
\begin{tabular}{@{}>{\raggedright\arraybackslash}p{22.0mm}>{\raggedright\arraybackslash}p{31.188mm}>{\raggedright\arraybackslash}p{31.188mm}>{\raggedright\arraybackslash}p{31.188mm}@{}}
\toprule
\textbf{Case} & \textbf{Masculine plural} & \textbf{Neuter plural} & \textbf{Feminine plural} \\
\midrule
Nominative & ūre & ūru & ūra \\
Genitive & ūrra & ūrra & ūrra \\
Dative & ūrum & ūrum & ūrum \\
Accusative & ūre & ūru & ūra \\
\bottomrule
\end{tabular}
\par\endgroup
% course-panel:T12-P02:end
% mdtable:T12:end
% mdparagraph:P082:begin
Instrumental singular, recognition only: masculine ūre; neuter ūre.
% mdparagraph:P082:end

% teaching: NEW-U03-PAR-eower-poss-strong
% mdtable:T13:begin
% course-table-style: 2026-10
% course-panel:T13-P01:begin
\par\begingroup
\fontsize{10.5}{13.2}\selectfont
\setlength{\LTleft}{0pt}\setlength{\LTright}{\fill}
\noindent
\begin{tabular}{@{}>{\raggedright\arraybackslash}p{22.0mm}>{\raggedright\arraybackslash}p{31.188mm}>{\raggedright\arraybackslash}p{31.188mm}>{\raggedright\arraybackslash}p{31.188mm}@{}}
\toprule
\textbf{Case} & \textbf{Masculine singular} & \textbf{Neuter singular} & \textbf{Feminine singular} \\
\midrule
Nominative & ēower & ēower & ēowru \\
Genitive & ēowres & ēowres & ēowerre \\
Dative & ēowrum & ēowrum & ēowerre \\
Accusative & ēowerne & ēower & ēowre \\
\bottomrule
\end{tabular}
\par\endgroup
% course-panel:T13-P01:end
\par\addvspace{5pt}
% course-panel:T13-P02:begin
\par\begingroup
\fontsize{10.5}{13.2}\selectfont
\setlength{\LTleft}{0pt}\setlength{\LTright}{\fill}
\noindent
\begin{tabular}{@{}>{\raggedright\arraybackslash}p{22.0mm}>{\raggedright\arraybackslash}p{31.188mm}>{\raggedright\arraybackslash}p{31.188mm}>{\raggedright\arraybackslash}p{31.188mm}@{}}
\toprule
\textbf{Case} & \textbf{Masculine plural} & \textbf{Neuter plural} & \textbf{Feminine plural} \\
\midrule
Nominative & ēowre & ēowru & ēowra \\
Genitive & ēowerra & ēowerra & ēowerra \\
Dative & ēowrum & ēowrum & ēowrum \\
Accusative & ēowre & ēowru & ēowra \\
\bottomrule
\end{tabular}
\par\endgroup
% course-panel:T13-P02:end
% mdtable:T13:end
% mdparagraph:P083:begin
Instrumental singular, recognition only: masculine ēowre; neuter ēowre.
% mdparagraph:P083:end

% mdparagraph:P084:begin
Ūre has feminine nominative singular ūre, as well as masculine and neuter nominative singular ūre; its selected neuter plural is ūru. Keep the table’s ūrre and ūrra distinct from the supported alternatives ūre and ūra. Ēower has ēowr- before vowel endings and ēower- before consonantal endings. Its supported unsyncopated -u alternative is ēoweru; single-r ēowere/\allowbreak{}ēowera also occur. These variants are valid alternatives, not reasons to ignore the selected paradigm when it is requested.
% mdparagraph:P084:end

\reserve{7\baselineskip}
% mdparagraph:P085:begin
\subsection*{Possessives — P1–P6}
% mdparagraph:P085:end

\par\begin{minipage}{\linewidth}\raggedright
% mdparagraph:P086:begin
\textbf{P1.}
% mdparagraph:P086:end

% item: NEW-U03-P01
% mdparagraph:P087:begin
Give the complete nominative, genitive, dative and accusative of the agreeing possessive mīn in all three genders, singular and plural. State what determines its ending.
% mdparagraph:P087:end

\end{minipage}\par\addvspace{4pt}

\par\begin{minipage}{\linewidth}\raggedright
% mdparagraph:P088:begin
\textbf{P2.}
% mdparagraph:P088:end

% item: NEW-U03-P02
% mdparagraph:P089:begin
Translate “I hear your bird. You hear my birds.” Address one person.
% mdparagraph:P089:end

\end{minipage}\par\addvspace{4pt}

\par\begin{minipage}{\linewidth}\raggedright
% mdparagraph:P090:begin
\textbf{P3.}
% mdparagraph:P090:end

% item: NEW-U03-P03
% mdparagraph:P091:begin
Translate “We two have a dwelling. Our dwelling is small. You two have a tree. Your tree is great.” Retain the two-person possessors.
% mdparagraph:P091:end

\end{minipage}\par\addvspace{4pt}

\par\begin{minipage}{\linewidth}\raggedright
% mdparagraph:P092:begin
\textbf{P4.}
% mdparagraph:P092:end

% item: NEW-U03-P04
% mdparagraph:P093:begin
Translate “Our birds live in your dwelling.” The possessors are several speakers and several addressees.
% mdparagraph:P093:end

\end{minipage}\par\addvspace{4pt}

\par\begin{minipage}{\linewidth}\raggedright
% mdparagraph:P094:begin
\textbf{P5.}
% mdparagraph:P094:end

% item: NEW-U03-P05
% mdparagraph:P095:begin
Translate “His birds live in my dwelling. Their dwelling is greater than our dwelling.” Use a male singular possessor for “his”.
% mdparagraph:P095:end

\end{minipage}\par\addvspace{4pt}

\par\begin{minipage}{\linewidth}\raggedright
% mdparagraph:P096:begin
\textbf{P6.}
% mdparagraph:P096:end

% item: NEW-U03-P06
% mdparagraph:P097:begin
The phrases mīne fuglas and uncer fugol are clause subjects; ēowrum hofe follows locative on. Identify each possessor’s person and number, then the possessed noun’s gender, number and case. Which information comes from the possessive stem, and which from its ending?
% mdparagraph:P097:end

\end{minipage}\par\addvspace{4pt}

\reserve{7\baselineskip}
% mdparagraph:P098:begin
\section*{7. Comparison}
% mdparagraph:P098:end

% teaching: NEW-U03-GRAM-comparison-regular
% teaching: NEW-U03-GRAM-comparison-umlaut
% teaching: NEW-U03-GRAM-comparison-irregular
% teaching: NEW-U03-VOC-leof-comp
% teaching: NEW-U03-VOC-leof-superl
% teaching: NEW-U03-PAR-leof-comp-weak
% teaching: NEW-U03-PAR-leof-superl-strong
% teaching: NEW-U03-PAR-leof-superl-weak
% teaching: NEW-U03-VOC-eald-comp
% teaching: NEW-U03-VOC-eald-superl
% teaching: NEW-U03-PAR-eald-comp-weak
% teaching: NEW-U03-PAR-eald-superl-strong
% teaching: NEW-U03-PAR-eald-superl-weak
% teaching: NEW-U03-VOC-god-adjective-comp
% teaching: NEW-U03-VOC-god-adjective-superl
% teaching: NEW-U03-PAR-god-adjective-comp-weak
% teaching: NEW-U03-PAR-god-adjective-superl-strong
% teaching: NEW-U03-PAR-god-adjective-superl-weak
% teaching: NEW-U03-VOC-yfel-comp
% teaching: NEW-U03-VOC-yfel-superl
% teaching: NEW-U03-PAR-yfel-comp-weak
% teaching: NEW-U03-PAR-yfel-superl-strong
% teaching: NEW-U03-PAR-yfel-superl-weak
% teaching: NEW-U03-VOC-micel-comp
% teaching: NEW-U03-VOC-micel-superl
% teaching: NEW-U03-PAR-micel-comp-weak
% teaching: NEW-U03-PAR-micel-superl-strong
% teaching: NEW-U03-PAR-micel-superl-weak
% teaching: NEW-U03-VOC-lytel-comp
% teaching: NEW-U03-VOC-lytel-superl
% teaching: NEW-U03-PAR-lytel-comp-weak
% teaching: NEW-U03-PAR-lytel-superl-strong
% teaching: NEW-U03-PAR-lytel-superl-weak
% mdparagraph:P099:begin
The positive names a quality; the comparative expresses a greater or lesser degree; the superlative expresses the highest or lowest degree. The regular series adds -ra and -ost. Eald has an umlaut series; the four remaining series have forms to learn separately. The comparative column gives masculine nominative singular. The superlative column gives the strong masculine/\allowbreak{}neuter nominative singular.
% mdparagraph:P099:end

% mdtable:T14:begin
\par\begingroup
\fontsize{10.5}{13.2}\selectfont
\setlength{\LTleft}{0pt}\setlength{\LTright}{\fill}
\noindent
\begin{tabular}{@{}>{\raggedright\arraybackslash}p{20.94mm}>{\raggedright\arraybackslash}p{23.03mm}>{\raggedright\arraybackslash}p{23.03mm}>{\raggedright\arraybackslash}p{48.57mm}@{}}
\toprule
\textbf{Positive} & \textbf{Comparative} & \textbf{Superlative} & \textbf{Degree meanings} \\
\midrule
lēof & lēofra & lēofost & dearer; dearest \\
eald & ieldra & ieldest & older; oldest \\
gōd & bętera & bętst & better; best \\
yfel & wiersa & wierrest & worse; worst \\
micel & māra & mǣst & greater/\allowbreak{}larger; greatest/\allowbreak{}largest \\
lȳtel & lǣssa & lǣst & smaller; smallest \\
\bottomrule
\end{tabular}
\par\endgroup

% mdtable:T14:end
% teaching: NEW-U03-GRAM-comparative-weak
\reserve{7\baselineskip}
% mdparagraph:P100:begin
\subsection*{Complete comparative inflection}
% mdparagraph:P100:end

% mdparagraph:P101:begin
A comparative is weak, including when it is a predicate with no preceding article. Māra supplies a complete model:
% mdparagraph:P101:end

% mdtable:T15:begin
% course-table-style: 2026-10
% course-panel:T15-P01:begin
\par\begingroup
\fontsize{10.5}{13.2}\selectfont
\setlength{\LTleft}{0pt}\setlength{\LTright}{\fill}
\noindent
\begin{tabular}{@{}>{\raggedright\arraybackslash}p{22.0mm}>{\raggedright\arraybackslash}p{31.188mm}>{\raggedright\arraybackslash}p{31.188mm}>{\raggedright\arraybackslash}p{31.188mm}@{}}
\toprule
\textbf{Case} & \textbf{Masculine singular} & \textbf{Neuter singular} & \textbf{Feminine singular} \\
\midrule
Nominative & māra & māre & māre \\
Genitive & māran & māran & māran \\
Dative & māran & māran & māran \\
Accusative & māran & māre & māran \\
\bottomrule
\end{tabular}
\par\endgroup
% course-panel:T15-P01:end
\par\addvspace{5pt}
% course-panel:T15-P02:begin
\par\begingroup
\fontsize{10.5}{13.2}\selectfont
\setlength{\LTleft}{0pt}\setlength{\LTright}{\fill}
\noindent
\begin{tabular}{@{}>{\raggedright\arraybackslash}p{22.0mm}>{\raggedright\arraybackslash}p{31.188mm}>{\raggedright\arraybackslash}p{31.188mm}>{\raggedright\arraybackslash}p{31.188mm}@{}}
\toprule
\textbf{Case} & \textbf{Masculine plural} & \textbf{Neuter plural} & \textbf{Feminine plural} \\
\midrule
Nominative & māran & māran & māran \\
Genitive & mārra & mārra & mārra \\
Dative & mārum & mārum & mārum \\
Accusative & māran & māran & māran \\
\bottomrule
\end{tabular}
\par\endgroup
% course-panel:T15-P02:end
% mdtable:T15:end
% mdparagraph:P102:begin
For each other comparative, replace its final masculine nominative -a with -e for neuter/\allowbreak{}feminine nominative singular and neuter accusative singular; with -an for the other singular cases and all nominative/\allowbreak{}accusative plurals; and with -um for all dative plurals. Use the following genitive plurals, which are the same in all genders. These explicit forms complete each comparative paradigm.
% mdparagraph:P102:end

% mdtable:T16:begin
\par\begingroup
\fontsize{10.5}{13.2}\selectfont
\setlength{\LTleft}{0pt}\setlength{\LTright}{\fill}
\noindent
\begin{tabular}{@{}>{\raggedright\arraybackslash}p{65.35mm}>{\raggedright\arraybackslash}p{55.84mm}@{}}
\toprule
\textbf{Masculine nom. sg.} & \textbf{Genitive plural} \\
\midrule
lēofra & lēofra \\
ieldra & ieldra \\
bętera & bęterra \\
wiersa & wiersena \\
māra & mārra \\
lǣssa & lǣssena \\
\bottomrule
\end{tabular}
\par\endgroup

% mdtable:T16:end
% mdparagraph:P103:begin
Do not form an extra sequence of r’s by blindly appending -ra to the comparative stem. The selected -ena in wiersena and lǣssena is a licensed weak ending.
% mdparagraph:P103:end

% teaching: NEW-U03-GRAM-superlative
\reserve{7\baselineskip}
% mdparagraph:P104:begin
\subsection*{Superlatives}
% mdparagraph:P104:end

% mdparagraph:P105:begin
Every listed superlative takes the complete strong and weak gōd patterns, replacing gōd- with the printed superlative stem. All six are selected long-stem patterns: strong feminine nominative singular and neuter N/\allowbreak{}A plural are endingless; strong instrumental singular has -e. An ordinary attributive superlative is normally weak; the taught predicative use is strong. The weak masculine nominatives are:
% mdparagraph:P105:end

% mdtable:T17:begin
\par\begingroup
\fontsize{10.5}{13.2}\selectfont
\setlength{\LTleft}{0pt}\setlength{\LTright}{\fill}
\noindent
\begin{tabular}{@{}>{\raggedright\arraybackslash}p{45.74mm}>{\raggedright\arraybackslash}p{34.41mm}>{\raggedright\arraybackslash}p{38.23mm}@{}}
\toprule
\textbf{Strong headword/\allowbreak{}stem} & \textbf{Weak m. nom. sg.} & \textbf{Weak n./\allowbreak{}f. nom. sg.} \\
\midrule
lēofost & lēofosta & lēofoste \\
ieldest & ieldesta & ieldeste \\
bętst & bętsta & bętste \\
wierrest & wierresta & wierreste \\
mǣst & mǣsta & mǣste \\
lǣst & lǣsta & lǣste \\
\bottomrule
\end{tabular}
\par\endgroup

% mdtable:T17:end
% mdparagraph:P106:begin
Use the same stem before all other strong/\allowbreak{}weak endings. This supplies complete inflection, not new partitive or substantival uses of the superlative.
% mdparagraph:P106:end

% teaching: NEW-U03-GRAM-comparison-thonne
% teaching: NEW-U03-VOC-thonne
\reserve{7\baselineskip}
% mdparagraph:P107:begin
\subsection*{A nominative standard with ðonne}
% mdparagraph:P107:end

% mdparagraph:P108:begin
Ðonne means “than” in the comparison taught here. When two subjects are compared with the copula, the second subject is nominative. The repeated copula can remain unexpressed. Restrict this pattern to the parallel copular comparison; it is not a rule for every clause containing “than”.
% mdparagraph:P108:end

% example: NEW-U03-EXAMPLE-07
% mdparagraph:P109:begin
Se mann is ieldra ðonne se fugol.
% mdparagraph:P109:end

% mdparagraph:P110:begin
The man is older than the bird.
% mdparagraph:P110:end

% mdparagraph:P111:begin
Mann and fugol are the contrasted nominative subjects. Ieldra is weak masculine nominative singular although it is a predicate.
% mdparagraph:P111:end

\reserve{7\baselineskip}
% mdparagraph:P112:begin
\subsection*{Comparison — C1–C7}
% mdparagraph:P112:end

\par\begin{minipage}{\linewidth}\raggedright
% mdparagraph:P113:begin
\textbf{C1.}
% mdparagraph:P113:end

% item: NEW-U03-C01
% mdparagraph:P114:begin
Give the selected positive, comparative and superlative masculine nominative singular of lēof, eald, gōd, micel, lȳtel and yfel. Use the weak superlative forms.
% mdparagraph:P114:end

\end{minipage}\par\addvspace{4pt}

\par\begin{minipage}{\linewidth}\raggedright
% mdparagraph:P115:begin
\textbf{C2.}
% mdparagraph:P115:end

% item: NEW-U03-C02
% mdparagraph:P116:begin
Give the complete weak nominative, genitive, dative and accusative of māra in all three genders, singular and plural.
% mdparagraph:P116:end

\end{minipage}\par\addvspace{4pt}

\par\begin{minipage}{\linewidth}\raggedright
% mdparagraph:P117:begin
\textbf{C3.}
% mdparagraph:P117:end

% item: NEW-U03-C03
% mdparagraph:P118:begin
Give the strong nominative singular of mǣst in all three genders and its weak nominative singular in all three genders. Then give the genitive and dative plurals in both declensions.
% mdparagraph:P118:end

\end{minipage}\par\addvspace{4pt}

\par\begin{minipage}{\linewidth}\raggedright
% mdparagraph:P119:begin
\textbf{C4.}
% mdparagraph:P119:end

% item: NEW-U03-C04
% mdparagraph:P120:begin
Translate “The great tree is greater than the small tree.” Identify the case of both trees and explain why the comparative has its selected ending.
% mdparagraph:P120:end

\end{minipage}\par\addvspace{4pt}

\par\begin{minipage}{\linewidth}\raggedright
% mdparagraph:P121:begin
\textbf{C5.}
% mdparagraph:P121:end

% item: NEW-U03-C05
% mdparagraph:P122:begin
Translate “The bird is dearer than the fish. The old dwelling is better than the great dwelling.”
% mdparagraph:P122:end

\end{minipage}\par\addvspace{4pt}

\par\begin{minipage}{\linewidth}\raggedright
% mdparagraph:P123:begin
\textbf{C6.}
% mdparagraph:P123:end

% item: NEW-U03-C06
% mdparagraph:P124:begin
Translate these nominative phrases: “the oldest tree”; “the dearest bird”; “the greatest work”; “the best man”; “the smallest star”.
% mdparagraph:P124:end

\end{minipage}\par\addvspace{4pt}

\par\begin{minipage}{\linewidth}\raggedright
% mdparagraph:P125:begin
\textbf{C7.}
% mdparagraph:P125:end

% item: NEW-U03-C07
% mdparagraph:P126:begin
Explain and correct the adjective ending in Se fugol is lǣss. Keep the meaning “The bird is smaller.” Then change the corrected sentence to refer to several birds.
% mdparagraph:P126:end

\end{minipage}\par\addvspace{4pt}

\reserve{7\baselineskip}
% mdparagraph:P127:begin
\section*{8. Adverbs and degree}
% mdparagraph:P127:end

% teaching: NEW-U03-GRAM-adverb-formation
% teaching: NEW-U03-GRAM-adverb-use
% teaching: NEW-U03-GRAM-adverb-degree
% teaching: NEW-U03-GRAM-adverb-comparison
% teaching: NEW-U03-VOC-wide
% teaching: NEW-U03-VOC-wel
% teaching: NEW-U03-VOC-swithe
% teaching: NEW-U03-VOC-sothlice
% teaching: NEW-U03-VOC-wid-citation
% teaching: NEW-U03-VOC-swith-citation
% teaching: NEW-U03-VOC-soth-citation
% teaching: NEW-U03-PAR-wide
% teaching: NEW-U03-PAR-wel
% teaching: NEW-U03-PAR-swithe
% teaching: NEW-U03-PAR-sothlice
% mdparagraph:P128:begin
An adverb qualifies an action, a statement, an adjective or another adverb. It does not agree in case, number or gender. The selected adjective citations wīd “wide”, swīð “strong” and sōð “true” are supplied bases for this formation; their adjective declensions are not assigned.
% mdparagraph:P128:end

% mdtable:T18:begin
\par\begingroup
\fontsize{10.5}{13.2}\selectfont
\setlength{\LTleft}{0pt}\setlength{\LTright}{\fill}
\noindent
\begin{tabular}{@{}>{\raggedright\arraybackslash}p{24.43mm}>{\raggedright\arraybackslash}p{28.15mm}>{\raggedright\arraybackslash}p{31.0mm}>{\raggedright\arraybackslash}p{31.99mm}@{}}
\toprule
\textbf{Positive adverb} & \textbf{Meaning} & \textbf{Comparative} & \textbf{Superlative} \\
\midrule
wīde & widely; far and wide & wīdor, more widely & wīdost, most widely \\
wel & well & bęt, better & bętst, best \\
swīðe & very & — & — \\
sōðlice & truly; indeed & — & — \\
\bottomrule
\end{tabular}
\par\endgroup

% mdtable:T18:end
% mdparagraph:P129:begin
The first derivation uses -e; sōðlice uses -lice. The regular adverbial comparative/\allowbreak{}superlative endings are -or/\allowbreak{}-ost. Memorize wel/\allowbreak{}bęt/\allowbreak{}bętst separately. An agreeing comparative adjective and the invariant adverb bęt perform different jobs. Swīðe can qualify an adjective. Sōðlice “truly” here is a new meaning distinct from the supplied connective use in the source reading.
% mdparagraph:P129:end

% example: NEW-U03-EXAMPLE-08
% mdparagraph:P130:begin
Hīe geworhton ðā hofu wel.
% mdparagraph:P130:end

% mdparagraph:P131:begin
They built the dwellings well.
% mdparagraph:P131:end

% mdparagraph:P132:begin
Wel qualifies geworhton; the dwellings remain an accusative object with their own morphology.
% mdparagraph:P132:end

\reserve{7\baselineskip}
% mdparagraph:P133:begin
\subsection*{Adverbs — V1–V4}
% mdparagraph:P133:end

\par\begin{minipage}{\linewidth}\raggedright
% mdparagraph:P134:begin
\textbf{V1.}
% mdparagraph:P134:end

% item: NEW-U03-V01
% mdparagraph:P135:begin
Using the supplied adjective citation wīd, give the positive, comparative and superlative adverbs. Also form the adverb meaning “truly” from the supplied citation sōð. State which of these adverbs changes for gender, case or number.
% mdparagraph:P135:end

\end{minipage}\par\addvspace{4pt}

\par\begin{minipage}{\linewidth}\raggedright
% mdparagraph:P136:begin
\textbf{V2.}
% mdparagraph:P136:end

% item: NEW-U03-V02
% mdparagraph:P137:begin
Translate “The men built the good dwelling well.” Identify the adjective and the adverb, and state what each modifies.
% mdparagraph:P137:end

\end{minipage}\par\addvspace{4pt}

\par\begin{minipage}{\linewidth}\raggedright
% mdparagraph:P138:begin
\textbf{V3.}
% mdparagraph:P138:end

% item: NEW-U03-V03
% mdparagraph:P139:begin
Translate “He built the dwelling better.” Use the adverb meaning “better”. Then replace it with the adverb meaning “best”. Explain why neither word agrees with the dwelling.
% mdparagraph:P139:end

\end{minipage}\par\addvspace{4pt}

\par\begin{minipage}{\linewidth}\raggedright
% mdparagraph:P140:begin
\textbf{V4.}
% mdparagraph:P140:end

% item: NEW-U03-V04
% mdparagraph:P141:begin
Translate Sōðlice, ðæt weorc is swīðe gōd. Compare it with Ðæt weorc is gōd. What do the two adverbs add, and why do they have no adjective agreement endings?
% mdparagraph:P141:end

\end{minipage}\par\addvspace{4pt}

\reserve{7\baselineskip}
% mdparagraph:P142:begin
\section*{9. Cardinal numbers: one through twelve}
% mdparagraph:P142:end

% teaching: NEW-U03-GRAM-cardinal-one
% teaching: NEW-U03-GRAM-cardinal-two-three
% teaching: NEW-U03-GRAM-cardinal-four-twelve
% teaching: NEW-U03-GRAM-cardinal-counting
% teaching: NEW-U03-VOC-an
% teaching: NEW-U03-VOC-twegen
% teaching: NEW-U03-VOC-thrie
% teaching: NEW-U03-VOC-feower
% teaching: NEW-U03-VOC-fif
% teaching: NEW-U03-VOC-siex
% teaching: NEW-U03-VOC-seofon
% teaching: NEW-U03-VOC-eahta
% teaching: NEW-U03-VOC-nigon
% teaching: NEW-U03-VOC-tien
% teaching: NEW-U03-VOC-endlefan
% teaching: NEW-U03-VOC-twelf
% mdparagraph:P143:begin
Ān modifies a singular count noun. Two and three take plural count nouns and inflect for the noun’s gender and case. Four through twelve are used here in their ordinary invariant position before plural count nouns. The noun keeps the case required by its clause; counting does not supply a new noun dual. A dual personal subject still takes a plural verb, as in Unit 2.
% mdparagraph:P143:end

\reserve{7\baselineskip}
% mdparagraph:P144:begin
\subsection*{Ān}
% mdparagraph:P144:end

% teaching: NEW-U03-PAR-an-strong
% mdtable:T19:begin
% course-table-style: 2026-10
% course-panel:T19-P01:begin
\par\begingroup
\fontsize{10.5}{13.2}\selectfont
\setlength{\LTleft}{0pt}\setlength{\LTright}{\fill}
\noindent
\begin{tabular}{@{}>{\raggedright\arraybackslash}p{22.0mm}>{\raggedright\arraybackslash}p{31.188mm}>{\raggedright\arraybackslash}p{31.188mm}>{\raggedright\arraybackslash}p{31.188mm}@{}}
\toprule
\textbf{Case} & \textbf{Masculine singular} & \textbf{Neuter singular} & \textbf{Feminine singular} \\
\midrule
Nominative & ān & ān & ān \\
Genitive & ānes & ānes & ānre \\
Dative & ānum & ānum & ānre \\
Accusative & ānne & ān & āne \\
\bottomrule
\end{tabular}
\par\endgroup
% course-panel:T19-P01:end
\par\addvspace{5pt}
% course-panel:T19-P02:begin
\par\begingroup
\fontsize{10.5}{13.2}\selectfont
\setlength{\LTleft}{0pt}\setlength{\LTright}{\fill}
\noindent
\begin{tabular}{@{}>{\raggedright\arraybackslash}p{22.0mm}>{\raggedright\arraybackslash}p{31.188mm}>{\raggedright\arraybackslash}p{31.188mm}>{\raggedright\arraybackslash}p{31.188mm}@{}}
\toprule
\textbf{Case} & \textbf{Masculine plural: reference only} & \textbf{Neuter plural: reference only} & \textbf{Feminine plural: reference only} \\
\midrule
Nominative & āne & ān & āna \\
Genitive & ānra & ānra & ānra \\
Dative & ānum & ānum & ānum \\
Accusative & āne & ān & āna \\
\bottomrule
\end{tabular}
\par\endgroup
% course-panel:T19-P02:end
% mdtable:T19:end
% mdparagraph:P145:begin
Instrumental singular, recognition only: masculine āne; neuter āne.
% mdparagraph:P145:end

% mdparagraph:P146:begin
Only the singular is used productively for cardinal “one”. The plural table is for recognition/\allowbreak{}reference and does not authorize “one” with plural count nouns. The alternative masculine accusative singular ǣnne is valid beside selected ānne. A weak form meaning “alone” and independent numeral uses are not assigned.
% mdparagraph:P146:end

\reserve{7\baselineskip}
% mdparagraph:P147:begin
\subsection*{Two and three}
% mdparagraph:P147:end

% teaching: NEW-U03-PAR-twegen
% mdparagraph:P148:begin
\textbf{twēgen}
% mdparagraph:P148:end

% mdtable:T20:begin
% course-table-style: 2026-10
% course-panel:T20-P01:begin
\par\begingroup
\fontsize{10.5}{13.2}\selectfont
\setlength{\LTleft}{0pt}\setlength{\LTright}{\fill}
\noindent
\begin{tabular}{@{}>{\raggedright\arraybackslash}p{22.0mm}>{\raggedright\arraybackslash}p{31.188mm}>{\raggedright\arraybackslash}p{31.188mm}>{\raggedright\arraybackslash}p{31.188mm}@{}}
\toprule
\textbf{Case} & \textbf{Masculine} & \textbf{Neuter} & \textbf{Feminine} \\
\midrule
Nominative & twēgen & twā & twā \\
Genitive & twēgra & twēgra & twēgra \\
Dative & twǣm & twǣm & twǣm \\
Accusative & twēgen & twā & twā \\
\bottomrule
\end{tabular}
\par\endgroup
% course-panel:T20-P01:end
% mdtable:T20:end
% teaching: NEW-U03-PAR-thrie
% mdparagraph:P149:begin
\textbf{ðrīe}
% mdparagraph:P149:end

% mdtable:T21:begin
% course-table-style: 2026-10
% course-panel:T21-P01:begin
\par\begingroup
\fontsize{10.5}{13.2}\selectfont
\setlength{\LTleft}{0pt}\setlength{\LTright}{\fill}
\noindent
\begin{tabular}{@{}>{\raggedright\arraybackslash}p{22.0mm}>{\raggedright\arraybackslash}p{31.188mm}>{\raggedright\arraybackslash}p{31.188mm}>{\raggedright\arraybackslash}p{31.188mm}@{}}
\toprule
\textbf{Case} & \textbf{Masculine} & \textbf{Neuter} & \textbf{Feminine} \\
\midrule
Nominative & ðrīe & ðrēo & ðrēo \\
Genitive & ðrēora & ðrēora & ðrēora \\
Dative & ðrīm & ðrīm & ðrīm \\
Accusative & ðrīe & ðrēo & ðrēo \\
\bottomrule
\end{tabular}
\par\endgroup
% course-panel:T21-P01:end
% mdtable:T21:end
% mdparagraph:P150:begin
All three genders share each numeral’s genitive and dative. Nominative and accusative plural distinguish the selected gender forms. A recipient still takes the dative numeral and dative noun.
% mdparagraph:P150:end

% example: NEW-U03-EXAMPLE-09
% mdparagraph:P151:begin
Wit habbað twā trēowu.
% mdparagraph:P151:end

% mdparagraph:P152:begin
We two have two trees.
% mdparagraph:P152:end

% mdparagraph:P153:begin
Wit identifies the speaker and one other; twā counts the neuter plural trees. These are separate grammatical choices.
% mdparagraph:P153:end

% teaching: NEW-U03-PAR-feower
% teaching: NEW-U03-PAR-fif
% teaching: NEW-U03-PAR-siex
% teaching: NEW-U03-PAR-seofon
% teaching: NEW-U03-PAR-eahta
% teaching: NEW-U03-PAR-nigon
% teaching: NEW-U03-PAR-tien
% teaching: NEW-U03-PAR-endlefan
% teaching: NEW-U03-PAR-twelf
% mdtable:T22:begin
\par\begingroup
\fontsize{10.5}{13.2}\selectfont
\setlength{\LTleft}{0pt}\setlength{\LTright}{\fill}
\noindent
\begin{tabular}{@{}>{\raggedright\arraybackslash}p{54.96mm}>{\raggedright\arraybackslash}p{66.23mm}@{}}
\toprule
\textbf{Number} & \textbf{Selected cardinal} \\
\midrule
4 & fēower \\
5 & fīf \\
6 & siex \\
7 & seofon \\
8 & eahta \\
9 & nigon \\
10 & tīen \\
11 & endlefan \\
12 & twęlf \\
\bottomrule
\end{tabular}
\par\endgroup

% mdtable:T22:end
% mdparagraph:P154:begin
Use these nine forms without case endings in the stated pre-nominal counting construction. Higher numbers, partitive genitives, and the inflection of postposed or substantival numerals are outside the present tasks.
% mdparagraph:P154:end

\reserve{7\baselineskip}
% mdparagraph:P155:begin
\subsection*{Cardinals — N1–N5}
% mdparagraph:P155:end

\par\begin{minipage}{\linewidth}\raggedright
% mdparagraph:P156:begin
\textbf{N1.}
% mdparagraph:P156:end

% item: NEW-U03-N01
% mdparagraph:P157:begin
Give the nominative, genitive, dative and accusative singular of ān in all three genders. Use the selected cardinal forms.
% mdparagraph:P157:end

\end{minipage}\par\addvspace{4pt}

\par\begin{minipage}{\linewidth}\raggedright
% mdparagraph:P158:begin
\textbf{N2.}
% mdparagraph:P158:end

% item: NEW-U03-N02
% mdparagraph:P159:begin
Give the nominative, genitive, dative and accusative of twēgen and ðrīe in all three genders. State which number their accompanying count nouns have.
% mdparagraph:P159:end

\end{minipage}\par\addvspace{4pt}

\par\begin{minipage}{\linewidth}\raggedright
% mdparagraph:P160:begin
\textbf{N3.}
% mdparagraph:P160:end

% item: NEW-U03-N03
% mdparagraph:P161:begin
Translate these nominative phrases: “two men”; “two herbs”; “two dwellings”; “three birds”; “three nights”; “three trees”.
% mdparagraph:P161:end

\end{minipage}\par\addvspace{4pt}

\par\begin{minipage}{\linewidth}\raggedright
% mdparagraph:P162:begin
\textbf{N4.}
% mdparagraph:P162:end

% item: NEW-U03-N04
% mdparagraph:P163:begin
Translate “We gave seed to two men. They gave water to three birds.”
% mdparagraph:P163:end

\end{minipage}\par\addvspace{4pt}

\par\begin{minipage}{\linewidth}\raggedright
% mdparagraph:P164:begin
\textbf{N5.}
% mdparagraph:P164:end

% item: NEW-U03-N05
% mdparagraph:P165:begin
Translate these nominative phrases: “four fish”; “five trees”; “six stars”; “seven days”; “eight birds”; “nine men”; “ten herbs”; “eleven years”; “twelve signs”. Explain the number of the nouns and the selected treatment of the cardinals.
% mdparagraph:P165:end

\end{minipage}\par\addvspace{4pt}

\reserve{7\baselineskip}
% mdparagraph:P166:begin
\section*{10. Ordinals: first through seventh}
% mdparagraph:P166:end

% teaching: NEW-U03-GRAM-ordinals
% teaching: NEW-U03-GRAM-ordinal-other
% teaching: NEW-U03-VOC-aeresta
% teaching: NEW-U03-VOC-other
% teaching: NEW-U03-VOC-thridda
% teaching: NEW-U03-VOC-feortha
% teaching: NEW-U03-VOC-fifta
% teaching: NEW-U03-VOC-siexta
% teaching: NEW-U03-VOC-seofotha
% mdtable:T23:begin
\par\begingroup
\fontsize{10.5}{13.2}\selectfont
\setlength{\LTleft}{0pt}\setlength{\LTright}{\fill}
\noindent
\begin{tabular}{@{}>{\raggedright\arraybackslash}p{24.69mm}>{\raggedright\arraybackslash}p{42.22mm}>{\raggedright\arraybackslash}p{51.47mm}@{}}
\toprule
\textbf{Order} & \textbf{Selected masculine nom. sg.} & \textbf{Declension/\allowbreak{}stem} \\
\midrule
1 & ǣresta & weak; ǣrest- \\
2 & ōðer & strong, special ōðer pattern \\
3 & ðridda & weak; ðridd- \\
4 & fēorða & weak; fēorð- \\
5 & fīfta & weak; fīft- \\
6 & siexta & weak; siext- \\
7 & seofoða & weak; seofoð- \\
\bottomrule
\end{tabular}
\par\endgroup

% mdtable:T23:end
% mdparagraph:P167:begin
Ordinals express position in an order. First and third through seventh are weak. The second, ōðer, remains strong even after the article. The case still follows the noun’s role; these ordinal forms do not introduce a new temporal use of on.
% mdparagraph:P167:end

\reserve{7\baselineskip}
% mdparagraph:P168:begin
\subsection*{Complete ordinary ordinal pattern}
% mdparagraph:P168:end

% teaching: NEW-U03-PAR-thridda-weak
% mdtable:T24:begin
% course-table-style: 2026-10
% course-panel:T24-P01:begin
\par\begingroup
\fontsize{10.5}{13.2}\selectfont
\setlength{\LTleft}{0pt}\setlength{\LTright}{\fill}
\noindent
\begin{tabular}{@{}>{\raggedright\arraybackslash}p{22.0mm}>{\raggedright\arraybackslash}p{31.188mm}>{\raggedright\arraybackslash}p{31.188mm}>{\raggedright\arraybackslash}p{31.188mm}@{}}
\toprule
\textbf{Case} & \textbf{Masculine singular} & \textbf{Neuter singular} & \textbf{Feminine singular} \\
\midrule
Nominative & ðridda & ðridde & ðridde \\
Genitive & ðriddan & ðriddan & ðriddan \\
Dative & ðriddan & ðriddan & ðriddan \\
Accusative & ðriddan & ðridde & ðriddan \\
\bottomrule
\end{tabular}
\par\endgroup
% course-panel:T24-P01:end
\par\addvspace{5pt}
% course-panel:T24-P02:begin
\par\begingroup
\fontsize{10.5}{13.2}\selectfont
\setlength{\LTleft}{0pt}\setlength{\LTright}{\fill}
\noindent
\begin{tabular}{@{}>{\raggedright\arraybackslash}p{22.0mm}>{\raggedright\arraybackslash}p{31.188mm}>{\raggedright\arraybackslash}p{31.188mm}>{\raggedright\arraybackslash}p{31.188mm}@{}}
\toprule
\textbf{Case} & \textbf{Masculine plural} & \textbf{Neuter plural} & \textbf{Feminine plural} \\
\midrule
Nominative & ðriddan & ðriddan & ðriddan \\
Genitive & ðridra & ðridra & ðridra \\
Dative & ðriddum & ðriddum & ðriddum \\
Accusative & ðriddan & ðriddan & ðriddan \\
\bottomrule
\end{tabular}
\par\endgroup
% course-panel:T24-P02:end
% mdtable:T24:end
% teaching: NEW-U03-PAR-aeresta-weak
% teaching: NEW-U03-PAR-feortha-weak
% teaching: NEW-U03-PAR-fifta-weak
% teaching: NEW-U03-PAR-siexta-weak
% teaching: NEW-U03-PAR-seofotha-weak
% mdparagraph:P169:begin
For first and fourth through seventh, remove masculine nominative -a to obtain the printed stem, then apply every weak gōd ending, including genitive plural -ra and dative plural -um. Ðridda has the exceptional genitive plural ðridra, with one d before r; retain dd in its other cells shown above. No strong alternative is substituted merely because an ordinal lacks an article.
% mdparagraph:P169:end

\reserve{7\baselineskip}
% mdparagraph:P170:begin
\subsection*{Ōðer}
% mdparagraph:P170:end

% teaching: NEW-U03-PAR-other-strong
% mdtable:T25:begin
% course-table-style: 2026-10
% course-panel:T25-P01:begin
\par\begingroup
\fontsize{10.5}{13.2}\selectfont
\setlength{\LTleft}{0pt}\setlength{\LTright}{\fill}
\noindent
\begin{tabular}{@{}>{\raggedright\arraybackslash}p{22.0mm}>{\raggedright\arraybackslash}p{31.188mm}>{\raggedright\arraybackslash}p{31.188mm}>{\raggedright\arraybackslash}p{31.188mm}@{}}
\toprule
\textbf{Case} & \textbf{Masculine singular} & \textbf{Neuter singular} & \textbf{Feminine singular} \\
\midrule
Nominative & ōðer & ōðer & ōðru \\
Genitive & ōðres & ōðres & ōðerre \\
Dative & ōðrum & ōðrum & ōðerre \\
Accusative & ōðerne & ōðer & ōðre \\
\bottomrule
\end{tabular}
\par\endgroup
% course-panel:T25-P01:end
\par\addvspace{5pt}
% course-panel:T25-P02:begin
\par\begingroup
\fontsize{10.5}{13.2}\selectfont
\setlength{\LTleft}{0pt}\setlength{\LTright}{\fill}
\noindent
\begin{tabular}{@{}>{\raggedright\arraybackslash}p{22.0mm}>{\raggedright\arraybackslash}p{31.188mm}>{\raggedright\arraybackslash}p{31.188mm}>{\raggedright\arraybackslash}p{31.188mm}@{}}
\toprule
\textbf{Case} & \textbf{Masculine plural} & \textbf{Neuter plural} & \textbf{Feminine plural} \\
\midrule
Nominative & ōðre & ōðru & ōðra \\
Genitive & ōðerra & ōðerra & ōðerra \\
Dative & ōðrum & ōðrum & ōðrum \\
Accusative & ōðre & ōðru & ōðra \\
\bottomrule
\end{tabular}
\par\endgroup
% course-panel:T25-P02:end
% mdtable:T25:end
% mdparagraph:P171:begin
Instrumental singular, recognition only: masculine ōðre; neuter ōðre.
% mdparagraph:P171:end

% mdparagraph:P172:begin
The selected feminine nominative singular is ōðru; ōðer and ōðeru are supported alternatives. Preserve ōðr- before the selected vowel endings, but ōðer- before -ne, -re and -ra. The article does not change this adjective to the weak pattern.
% mdparagraph:P172:end

% example: NEW-U03-EXAMPLE-10
% mdparagraph:P173:begin
Ūre mann is on ðǣm ōðrum hofe.
% mdparagraph:P173:end

% mdparagraph:P174:begin
Our man is in the second dwelling.
% mdparagraph:P174:end

% mdparagraph:P175:begin
Ōðrum is strong neuter dative singular after the article; hofe is dative after locative on.
% mdparagraph:P175:end

\reserve{7\baselineskip}
% mdparagraph:P176:begin
\subsection*{Ordinals — O1–O3}
% mdparagraph:P176:end

\par\begin{minipage}{\linewidth}\raggedright
% mdparagraph:P177:begin
\textbf{O1.}
% mdparagraph:P177:end

% item: NEW-U03-O01
% mdparagraph:P178:begin
Give the nominative and genitive singular of “the second day”, “the second night” and “the second year”. Explain the declension of ōðer after the article.
% mdparagraph:P178:end

\end{minipage}\par\addvspace{4pt}

\par\begin{minipage}{\linewidth}\raggedright
% mdparagraph:P179:begin
\textbf{O2.}
% mdparagraph:P179:end

% item: NEW-U03-O02
% mdparagraph:P180:begin
Give the complete weak nominative, genitive, dative and accusative of ðridda in all three genders, singular and plural.
% mdparagraph:P180:end

\end{minipage}\par\addvspace{4pt}

\par\begin{minipage}{\linewidth}\raggedright
% mdparagraph:P181:begin
\textbf{O3.}
% mdparagraph:P181:end

% item: NEW-U03-O03
% mdparagraph:P182:begin
Translate “Three birds live in the first dwelling. The old man gave food to the fifth bird.”
% mdparagraph:P182:end

\end{minipage}\par\addvspace{4pt}

\reserve{7\baselineskip}
% mdparagraph:P183:begin
\section*{11. Demonstratives}
% mdparagraph:P183:end

% teaching: NEW-U03-GRAM-near-demonstrative
% teaching: NEW-U03-GRAM-far-demonstrative
% teaching: NEW-U03-VOC-se
% teaching: NEW-U03-VOC-thes
% mdparagraph:P184:begin
The article forms se, sēo, ðæt can also point to “that” or “those”. Their forms are familiar, but this pointing meaning is now assigned. Ðes, ðēos, ðis means “this”; its plural means “these”. A demonstrative agrees with its head noun. An ordinary adjective following it is weak. Relative uses and general third-person pronoun uses of these words are not introduced.
% mdparagraph:P184:end

\reserve{7\baselineskip}
% mdparagraph:P185:begin
\subsection*{Se: that/\allowbreak{}those; also the familiar article}
% mdparagraph:P185:end

% teaching: NEW-U03-PAR-se-demonstrative
% mdtable:T26:begin
% course-table-style: 2026-10
% course-panel:T26-P01:begin
\par\begingroup
\fontsize{10.5}{13.2}\selectfont
\setlength{\LTleft}{0pt}\setlength{\LTright}{\fill}
\noindent
\begin{tabular}{@{}>{\raggedright\arraybackslash}p{22.0mm}>{\raggedright\arraybackslash}p{31.188mm}>{\raggedright\arraybackslash}p{31.188mm}>{\raggedright\arraybackslash}p{31.188mm}@{}}
\toprule
\textbf{Case} & \textbf{Masculine singular} & \textbf{Neuter singular} & \textbf{Feminine singular} \\
\midrule
Nominative & se & ðæt & sēo \\
Genitive & ðæs & ðæs & ðǣre \\
Dative & ðǣm & ðǣm & ðǣre \\
Accusative & ðone & ðæt & ðā \\
\bottomrule
\end{tabular}
\par\endgroup
% course-panel:T26-P01:end
\par\addvspace{5pt}
% course-panel:T26-P02:begin
\par\begingroup
\fontsize{10.5}{13.2}\selectfont
\setlength{\LTleft}{0pt}\setlength{\LTright}{\fill}
\noindent
\begin{tabular}{@{}>{\raggedright\arraybackslash}p{22.0mm}>{\raggedright\arraybackslash}p{31.188mm}>{\raggedright\arraybackslash}p{31.188mm}>{\raggedright\arraybackslash}p{31.188mm}@{}}
\toprule
\textbf{Case} & \textbf{Masculine plural} & \textbf{Neuter plural} & \textbf{Feminine plural} \\
\midrule
Nominative & ðā & ðā & ðā \\
Genitive & ðāra & ðāra & ðāra \\
Dative & ðǣm & ðǣm & ðǣm \\
Accusative & ðā & ðā & ðā \\
\bottomrule
\end{tabular}
\par\endgroup
% course-panel:T26-P02:end
% mdtable:T26:end
% mdparagraph:P186:begin
Instrumental singular, recognition only: masculine ðȳ; neuter ðȳ.
% mdparagraph:P186:end

% mdparagraph:P187:begin
The table repeats the existing article morphology. The instrumental alternatives ðȳ and ðon remain recognition forms.
% mdparagraph:P187:end

\reserve{7\baselineskip}
% mdparagraph:P188:begin
\subsection*{Ðes: this/\allowbreak{}these}
% mdparagraph:P188:end

% teaching: NEW-U03-PAR-thes-demonstrative
% mdtable:T27:begin
% course-table-style: 2026-10
% course-panel:T27-P01:begin
\par\begingroup
\fontsize{10.5}{13.2}\selectfont
\setlength{\LTleft}{0pt}\setlength{\LTright}{\fill}
\noindent
\begin{tabular}{@{}>{\raggedright\arraybackslash}p{22.0mm}>{\raggedright\arraybackslash}p{31.188mm}>{\raggedright\arraybackslash}p{31.188mm}>{\raggedright\arraybackslash}p{31.188mm}@{}}
\toprule
\textbf{Case} & \textbf{Masculine singular} & \textbf{Neuter singular} & \textbf{Feminine singular} \\
\midrule
Nominative & ðes & ðis & ðēos \\
Genitive & ðisses & ðisses & ðisse \\
Dative & ðissum & ðissum & ðisse \\
Accusative & ðisne & ðis & ðās \\
\bottomrule
\end{tabular}
\par\endgroup
% course-panel:T27-P01:end
\par\addvspace{5pt}
% course-panel:T27-P02:begin
\par\begingroup
\fontsize{10.5}{13.2}\selectfont
\setlength{\LTleft}{0pt}\setlength{\LTright}{\fill}
\noindent
\begin{tabular}{@{}>{\raggedright\arraybackslash}p{22.0mm}>{\raggedright\arraybackslash}p{31.188mm}>{\raggedright\arraybackslash}p{31.188mm}>{\raggedright\arraybackslash}p{31.188mm}@{}}
\toprule
\textbf{Case} & \textbf{Masculine plural} & \textbf{Neuter plural} & \textbf{Feminine plural} \\
\midrule
Nominative & ðās & ðās & ðās \\
Genitive & ðissa & ðissa & ðissa \\
Dative & ðissum & ðissum & ðissum \\
Accusative & ðās & ðās & ðās \\
\bottomrule
\end{tabular}
\par\endgroup
% course-panel:T27-P02:end
% mdtable:T27:end
% mdparagraph:P189:begin
Instrumental singular, recognition only: masculine ðȳs; neuter ðȳs.
% mdparagraph:P189:end

% mdparagraph:P190:begin
Keep Cook’s short-e spelling ðes. Do not choose a form from the English word alone: a recipient may require a dative singular demonstrative, while an object may require an accusative plural one.
% mdparagraph:P190:end

% example: NEW-U03-EXAMPLE-06
% mdparagraph:P191:begin
Wē lufiað ðis lȳtle dēor.
% mdparagraph:P191:end

% mdparagraph:P192:begin
We love this small animal.
% mdparagraph:P192:end

% mdparagraph:P193:begin
Ðis and lȳtle are neuter accusative singular with the object dēor; the adjective is weak after the demonstrative.
% mdparagraph:P193:end

\reserve{7\baselineskip}
% mdparagraph:P194:begin
\subsection*{Demonstratives — D1–D4}
% mdparagraph:P194:end

\par\begin{minipage}{\linewidth}\raggedright
% mdparagraph:P195:begin
\textbf{D1.}
% mdparagraph:P195:end

% item: NEW-U03-D01
% mdparagraph:P196:begin
Give the complete nominative, genitive, dative and accusative of ðes in all three singular genders and in the plural.
% mdparagraph:P196:end

\end{minipage}\par\addvspace{4pt}

\par\begin{minipage}{\linewidth}\raggedright
% mdparagraph:P197:begin
\textbf{D2.}
% mdparagraph:P197:end

% item: NEW-U03-D02
% mdparagraph:P198:begin
Translate “This tree is great. That tree is small. These birds live in that dwelling.”
% mdparagraph:P198:end

\end{minipage}\par\addvspace{4pt}

\par\begin{minipage}{\linewidth}\raggedright
% mdparagraph:P199:begin
\textbf{D3.}
% mdparagraph:P199:end

% item: NEW-U03-D03
% mdparagraph:P200:begin
Translate “I give seed to this man. I hear these birds.”
% mdparagraph:P200:end

\end{minipage}\par\addvspace{4pt}

\par\begin{minipage}{\linewidth}\raggedright
% mdparagraph:P201:begin
\textbf{D4.}
% mdparagraph:P201:end

% item: NEW-U03-D04
% mdparagraph:P202:begin
Translate these nominative phrases: “this good man”; “these great trees”; “good men”. Explain the declension of each adjective.
% mdparagraph:P202:end

\end{minipage}\par\addvspace{4pt}

\reserve{7\baselineskip}
% mdparagraph:P203:begin
\section*{12. Cook, Group I, Exercise 4}
% mdparagraph:P203:end

% mdparagraph:P204:begin
Translate items 1–17. For a noun-phrase fragment, use a nominative head unless a task explicitly gives another case. Retain the dependent genitive’s article and place that dependent before its head, as in the earlier Cook exercises. Complete sentences require the case appropriate to each clause role. Use the selected forms when practising a paradigm; the key distinguishes those targets from defensible alternatives.
% mdparagraph:P204:end

% mdparagraph:P205:begin
Cook’s introductory note reads: “[NOTE. — For adjectives preceded by a demonstrative, see 55.]” The reference is to his §55. Item 12 retains Cook’s literal reference “(80)”; it points to the ordinal discussion.
% mdparagraph:P205:end

\par\begin{minipage}{\linewidth}\raggedright
% item: EX-I-04-01
% mdparagraph:P206:begin
\itempara{1}{The great whales.}
% mdparagraph:P206:end

\end{minipage}\par\addvspace{4pt}

\par\begin{minipage}{\linewidth}\raggedright
% item: EX-I-04-02
% mdparagraph:P207:begin
\itempara{2}{The earth was empty and void.}
% mdparagraph:P207:end

\end{minipage}\par\addvspace{4pt}

\par\begin{minipage}{\linewidth}\raggedright
% item: EX-I-04-03
% mdparagraph:P208:begin
\itempara{3}{The things were good.}
% mdparagraph:P208:end

\end{minipage}\par\addvspace{4pt}

\par\begin{minipage}{\linewidth}\raggedright
% item: EX-I-04-04
% mdparagraph:P209:begin
\itempara{4}{The fruit-bearing trees.}
% mdparagraph:P209:end

\end{minipage}\par\addvspace{4pt}

\par\begin{minipage}{\linewidth}\raggedright
% item: EX-I-04-05
% mdparagraph:P210:begin
\itempara{5}{The great lights.}
% mdparagraph:P210:end

\end{minipage}\par\addvspace{4pt}

\par\begin{minipage}{\linewidth}\raggedright
% item: EX-I-04-06
% mdparagraph:P211:begin
\itempara{6}{The little star.}
% mdparagraph:P211:end

\end{minipage}\par\addvspace{4pt}

\par\begin{minipage}{\linewidth}\raggedright
% item: EX-I-04-07
% mdparagraph:P212:begin
\itempara{7}{The smaller bird.}
% mdparagraph:P212:end

\end{minipage}\par\addvspace{4pt}

\par\begin{minipage}{\linewidth}\raggedright
% item: EX-I-04-08
% mdparagraph:P213:begin
\itempara{8}{The larger tree.}
% mdparagraph:P213:end

\end{minipage}\par\addvspace{4pt}

\par\begin{minipage}{\linewidth}\raggedright
% item: EX-I-04-09
% mdparagraph:P214:begin
\itempara{9}{The life of the good man.}
% mdparagraph:P214:end

\end{minipage}\par\addvspace{4pt}

\par\begin{minipage}{\linewidth}\raggedright
% item: EX-I-04-10
% mdparagraph:P215:begin
\itempara{10}{The beginning of the seventh night.}
% mdparagraph:P215:end

\end{minipage}\par\addvspace{4pt}

\par\begin{minipage}{\linewidth}\raggedright
% item: EX-I-04-11
% mdparagraph:P216:begin
\itempara{11}{The waters of the great sea.}
% mdparagraph:P216:end

\end{minipage}\par\addvspace{4pt}

\par\begin{minipage}{\linewidth}\raggedright
% item: EX-I-04-12
% mdparagraph:P217:begin
\itempara{12}{The second (80) day and the second night.}
% mdparagraph:P217:end

\end{minipage}\par\addvspace{4pt}

\par\begin{minipage}{\linewidth}\raggedright
% item: EX-I-04-13
% mdparagraph:P218:begin
\itempara{13}{The third day of the second year.}
% mdparagraph:P218:end

\end{minipage}\par\addvspace{4pt}

\par\begin{minipage}{\linewidth}\raggedright
% item: EX-I-04-14
% mdparagraph:P219:begin
\itempara{14}{The food of the fifth man.}
% mdparagraph:P219:end

\end{minipage}\par\addvspace{4pt}

\par\begin{minipage}{\linewidth}\raggedright
% item: EX-I-04-15
% mdparagraph:P220:begin
\itempara{15}{The light of the sixth evening.}
% mdparagraph:P220:end

\end{minipage}\par\addvspace{4pt}

\par\begin{minipage}{\linewidth}\raggedright
% item: EX-I-04-16
% mdparagraph:P221:begin
\itempara{16}{The works of the moving creatures.}
% mdparagraph:P221:end

\end{minipage}\par\addvspace{4pt}

\par\begin{minipage}{\linewidth}\raggedright
% item: EX-I-04-17
% mdparagraph:P222:begin
\itempara{17}{The morning of the first day.}
% mdparagraph:P222:end

% deferred: EX-I-04-18
% mdparagraph:P223:begin
Cook I.4.18, “Increase (sing.) and be multiplied.”, is reserved for Unit 7, where its command and passive construction can be taught. It is not an additional task in this unit.
% mdparagraph:P223:end

\end{minipage}\par\addvspace{4pt}

\reserve{7\baselineskip}
% mdparagraph:P224:begin
\section*{13. Controlled reading: Two dwellings and three birds}
% mdparagraph:P224:end

% mdparagraph:P225:begin
This connected account is newly composed teaching material. The old man is one of the builders. In his closing speech, ‘our’ includes all the builders.
% mdparagraph:P225:end

% composition: NEW-U03-TEXT-DWELLINGS-01
% mdparagraph:P226:begin
\par\begin{minipage}{\linewidth}\raggedright
\readingtext{Ðā gōdan męnn geworhton twā hofu. Ān hof wæs micel, and ān hof wæs lȳtel. Ðæt micle hof wæs māre ðonne ðæt lȳtle hof. Hīe sealdon ðæt micle hof ðǣm ealdan męnn. Hē hæfde ðrīe fuglas. Hē sealde his fuglum sǣd and wæter. Ðā fuglas lifdon on ðǣm miclan hofe. Se mann sægde: “Mīne fuglas sindon gōde. Ūre weorc is gōd.”}
\end{minipage}\par
% mdparagraph:P226:end

% mdparagraph:P227:begin
All words and constructions in this account have been taught. Track the case of each noun phrase before deciding who builds, gives, possesses or lives somewhere. Inside the direct quotation, the speaker’s first-person possessives refer to him and to the group identified in the orientation.
% mdparagraph:P227:end

\reserve{7\baselineskip}
% mdparagraph:P228:begin
\subsection*{Controlled reading — L1–L5}
% mdparagraph:P228:end

\par\begin{minipage}{\linewidth}\raggedright
% mdparagraph:P229:begin
\textbf{L1.}
% mdparagraph:P229:end

% item: NEW-U03-L01
% mdparagraph:P230:begin
Translate the complete new account “Two dwellings and three birds”. Keep the givers, recipient and possessed things distinct.
% mdparagraph:P230:end

\end{minipage}\par\addvspace{4pt}

\par\begin{minipage}{\linewidth}\raggedright
% mdparagraph:P231:begin
\textbf{L2.}
% mdparagraph:P231:end

% item: NEW-U03-L02
% mdparagraph:P232:begin
List the things counted by twā, ān and ðrīe in the account. Identify the gender, number and case of each numeral in its clause.
% mdparagraph:P232:end

\end{minipage}\par\addvspace{4pt}

\par\begin{minipage}{\linewidth}\raggedright
% mdparagraph:P233:begin
\textbf{L3.}
% mdparagraph:P233:end

% item: NEW-U03-L03
% mdparagraph:P234:begin
In the closing quotation, identify the possessor and the possessed noun in Mīne fuglas and Ūre weorc. Explain why the same final -e occurs with different gender and number analyses.
% mdparagraph:P234:end

\end{minipage}\par\addvspace{4pt}

\par\begin{minipage}{\linewidth}\raggedright
% mdparagraph:P235:begin
\textbf{L4.}
% mdparagraph:P235:end

% item: NEW-U03-L04
% mdparagraph:P236:begin
In the account, identify the strong predicate adjectives, the weak attributive adjectives and the comparative. Give the noun or pronoun governing each agreement.
% mdparagraph:P236:end

\end{minipage}\par\addvspace{4pt}

\par\begin{minipage}{\linewidth}\raggedright
% mdparagraph:P237:begin
\textbf{L5.}
% mdparagraph:P237:end

% item: NEW-U03-L05
% mdparagraph:P238:begin
Retell the first four sentences of the account in the present. Preserve all people, things, cases and adjective endings.
% mdparagraph:P238:end

\end{minipage}\par\addvspace{4pt}

\reserve{7\baselineskip}
% mdparagraph:P239:begin
\section*{14. Source reading: the fourth day}
% mdparagraph:P239:end

% mdparagraph:P240:begin
This is the complete fourth-day account from Cook’s first reading. Read the local help before translating. Its wishes, purpose clauses and resultative formulae are supplied; the adjective, numeral, noun and familiar indicative forms can be analysed with the systems learned so far. Do not infer that every form of a familiar verb has been taught.
% mdparagraph:P240:end

% source: FB-I-CREATION-FOURTH-DAY-01
% mdparagraph:P241:begin
\par\begin{minipage}{\linewidth}\raggedright
\readingtext{God cwæð ðā sōðlice, “Bēon nū lēoht on ðǣre heofonan fæstnesse, and tōdǣlen dæg and niht, and bēon tō tācnum, and tō tīdum, and tō dagum, and tō gēarum. And hīe scīnen on ðǣre heofonan fæstnesse, and ālīehten ðā eorðan”; hit wæs ðā swā geworden. And God geworhte twā miclu lēoht: ðæt māre lēoht tō ðæs dæges līehtinge, and ðæt lǣsse lēoht tō ðǣre niht līehtinge; and steorran hē geworhte. And gesętte hīe on ðǣre heofonan, ðæt hīe scinen ofer eorðan, and gīemden ðæs dæges and ðǣre niht, and tōdǣlden lēoht and ðīestru; God geseah ðā ðæt hit gōd wæs. And wæs geworden ǣfen and męrgen se fēorða dæg.}
\end{minipage}\par
% mdparagraph:P241:end

\reserve{7\baselineskip}
% mdparagraph:P242:begin
\subsection*{Local forms and meanings}
% mdparagraph:P242:end

% teaching: NEW-U03-SUPPORT-cwethan
% teaching: NEW-U03-SUPPORT-tha-adverb
% teaching: NEW-U03-SUPPORT-sothlice
% teaching: NEW-U03-SUPPORT-wesan
% teaching: NEW-U03-SUPPORT-nu
% teaching: NEW-U03-SUPPORT-heofone
% teaching: NEW-U03-SUPPORT-todaelan
% teaching: NEW-U03-SUPPORT-to
% teaching: NEW-U03-SUPPORT-scinan
% teaching: NEW-U03-SUPPORT-aliehtan
% teaching: NEW-U03-SUPPORT-swa
% teaching: NEW-U03-SUPPORT-geweorthan
% teaching: NEW-U03-SUPPORT-gesettan
% teaching: NEW-U03-SUPPORT-thaet-conjunction
% teaching: NEW-U03-SUPPORT-ofer
% teaching: NEW-U03-SUPPORT-gieman
% teaching: NEW-U03-SUPPORT-geseon
% teaching: NEW-U03-SUPPORT-morgen
% mdtable:T28:begin
\par\begingroup
\fontsize{10.5}{13.2}\selectfont
\setlength{\LTleft}{0pt}\setlength{\LTright}{\fill}
\noindent
\begin{tabular}{@{}>{\raggedright\arraybackslash}p{29.38mm}>{\raggedright\arraybackslash}p{91.81mm}@{}}
\toprule
\textbf{Source expression} & \textbf{Local help} \\
\midrule
cwæð & said; supplied third-person singular past of cweðan. The productive sęcgan paradigm does not supply this form. \\
ðā; sōðlice; nū & then; connective sōðlice (Cook compares Latin autem); now. The connective is distinct from the adverb “truly” learned above. \\
Bēon /\allowbreak{} bēon & let there be /\allowbreak{} let them be: plural present optative forms, expressing a wish or command. They are not infinitives or the indicative sindon. \\
heofonan & weak feminine heofone “heaven”: genitive singular before fæstnesse, but dative singular in on ðǣre heofonan. This is not the strong masculine heofon paradigm. \\
tōdǣlen; scīnen; ālīehten & let them separate; let them shine; let them illuminate. Supplied plural present optatives with the lights as subject. \\
tō tācnum /\allowbreak{} tīdum /\allowbreak{} dagum /\allowbreak{} gēarum & for signs /\allowbreak{} seasons /\allowbreak{} days /\allowbreak{} years. Tō governs these datives in a supplied function or purpose construction. \\
swā; geworden & thus, so; come to pass/\allowbreak{}become, a supplied past participle. Translate the complete formula as given below. \\
tō … līehtinge & for the illumination of … ; dative līehtinge contains a dependent genitive day or night. \\
gesętte & he set/\allowbreak{}placed; supplied past form of the prefixed verb. Its subject continues God/\allowbreak{}hē, although no subject is repeated in this clause. Hīe is the accusative object “them”. \\
ðæt hīe scinen … & so that they might shine … . Here ðæt introduces purpose. Preserve the short i in past optative scinen, contrasting present scīnen earlier. \\
ofer eorðan & over the earth. The whole phrase is supplied. Eorðan can be accusative or dative singular in this spelling; no unique choice is required. \\
gīemden; tōdǣlden & might rule over; might separate. Supplied plural past optatives continuing the purpose clause. Gīemden governs the following genitives. \\
geseah ðā ðæt … & then saw that … ; supplied strong past verb and content clause. Here ðæt is a conjunction, not an article. \\
męrgen & source variant of the already learned morgen “morning”; retain its spelling when reading or quoting this passage. \\
\bottomrule
\end{tabular}
\par\endgroup

% mdtable:T28:end
\reserve{7\baselineskip}
% mdparagraph:P243:begin
\subsection*{Whole constructions}
% mdparagraph:P243:end

% teaching: NEW-U03-SUPPORT-U03-R-QUOTE
% mdparagraph:P244:begin
\textbf{Opening speech.} God cwæð ðā sōðlice = God then said; sōðlice is connective here (Latin autem).
% mdparagraph:P244:end

% teaching: NEW-U03-SUPPORT-U03-R-WISH
% mdparagraph:P245:begin
\textbf{The wish.} Bēon/\allowbreak{}bēon = let there be/\allowbreak{}let them be; tōdǣlen = let them separate; scīnen = let them shine; ālīehten = let them illuminate. All four verbs express the wish and take the lights as their plural subject.
% mdparagraph:P245:end

% teaching: NEW-U03-SUPPORT-U03-R-HEAVEN-GEN
% mdparagraph:P246:begin
\textbf{Nested heaven phrase.} on ðǣre heofonan fæstnesse = in the firmament of heaven. Heofonan is genitive singular and fæstnesse dative singular; ðǣre agrees with heofonan.
% mdparagraph:P246:end

% teaching: NEW-U03-SUPPORT-U03-R-TO-FUNCTION
% mdparagraph:P247:begin
\textbf{Assigned functions.} tō tācnum/\allowbreak{}tīdum/\allowbreak{}dagum/\allowbreak{}gēarum = for signs/\allowbreak{}seasons/\allowbreak{}days/\allowbreak{}years; tō ... līehtinge = for the illumination of ...; genitive day/\allowbreak{}night is nested inside the dative phrase.
% mdparagraph:P247:end

% teaching: NEW-U03-SUPPORT-U03-R-FULFILMENT
% mdparagraph:P248:begin
\textbf{Fulfilment.} hit wæs ðā swā geworden = it then came to pass thus /\allowbreak{} it was then so; geworden is supplied.
% mdparagraph:P248:end

% teaching: NEW-U03-SUPPORT-U03-R-SET-HEAVEN
% mdparagraph:P249:begin
\textbf{Setting the lights.} And gesętte hīe = and he set them. The subject continues God/\allowbreak{}hē. on ðǣre heofonan = in heaven, with source weak-feminine dative.
% mdparagraph:P249:end

% teaching: NEW-U03-SUPPORT-U03-R-PURPOSE
% mdparagraph:P250:begin
\textbf{Purpose clause.} ðæt hīe scinen ... and gīemden ... and tōdǣlden ... = so that they might shine ... and rule over ... and separate ... . Scinen has short i; the present scīnen occurs earlier in the speech.
% mdparagraph:P250:end

% teaching: NEW-U03-SUPPORT-U03-R-OVER
% mdparagraph:P251:begin
\textbf{Over the earth.} Ofer eorðan means “over the earth”. Either accusative or dative singular is possible here; the form eorðan does not decide between them.
% mdparagraph:P251:end

% teaching: NEW-U03-SUPPORT-U03-R-RULE-GEN
% mdparagraph:P252:begin
\textbf{Governed genitives.} gīemden ðæs dæges and ðǣre niht = might rule over the day and the night. Here the verb governs the two genitives; this is not an of-relationship between nouns.
% mdparagraph:P252:end

% teaching: NEW-U03-SUPPORT-U03-R-SEE-CONTENT
% mdparagraph:P253:begin
\textbf{Seeing that it was good.} God geseah ðā ðæt hit gōd wæs = God then saw that it was good. Ðæt is a conjunction, not the article. This clause pattern is supplied for the reading.
% mdparagraph:P253:end

% teaching: NEW-U03-SUPPORT-U03-R-CLOSE
% mdparagraph:P254:begin
\textbf{Closing day formula.} And wæs geworden ǣfen and męrgen se fēorða dæg = and evening and morning came to pass, the fourth day. This is a supplied formula, including its word order and singular wæs.
% mdparagraph:P254:end

% mdparagraph:P255:begin
The two occurrences of geworhte belong to the familiar indicative paradigm of gewyrcean. The paired greater and lesser lights use the new counted sense of lēoht. The text assigns their functions without naming the sun and moon. The final ordinal identifies the fourth day. The known wæs forms within fulfilment formulae do not make the larger constructions independent production material.
% mdparagraph:P255:end

\reserve{7\baselineskip}
% mdparagraph:P256:begin
\subsection*{Source reading — R1–R6}
% mdparagraph:P256:end

\par\begin{minipage}{\linewidth}\raggedright
% mdparagraph:P257:begin
\textbf{R1.}
% mdparagraph:P257:end

% item: NEW-U03-R01
% mdparagraph:P258:begin
Translate the complete fourth-day Creation account using the local support. Preserve the command, fulfilment, purposes and closing day formula.
% mdparagraph:P258:end

\end{minipage}\par\addvspace{4pt}

\par\begin{minipage}{\linewidth}\raggedright
% mdparagraph:P259:begin
\textbf{R2.}
% mdparagraph:P259:end

% item: NEW-U03-R02
% mdparagraph:P260:begin
In twā miclu lēoht, identify the gender, number and case of all three words. Give their headwords and explain the selected adjective stem and ending.
% mdparagraph:P260:end

\end{minipage}\par\addvspace{4pt}

\par\begin{minipage}{\linewidth}\raggedright
% mdparagraph:P261:begin
\textbf{R3.}
% mdparagraph:P261:end

% item: NEW-U03-R03
% mdparagraph:P262:begin
Compare ðæt māre lēoht and ðæt lǣsse lēoht. Give the positive adjectives and the analysis of each comparative. Using the local purpose support, state the function assigned to each light.
% mdparagraph:P262:end

\end{minipage}\par\addvspace{4pt}

\par\begin{minipage}{\linewidth}\raggedright
% mdparagraph:P263:begin
\textbf{R4.}
% mdparagraph:P263:end

% item: NEW-U03-R04
% mdparagraph:P264:begin
Find the two ðæt forms that introduce the light noun phrases. Contrast them with the ðæt before hīe scinen and before hit gōd wæs, using the supplied clause notes.
% mdparagraph:P264:end

\end{minipage}\par\addvspace{4pt}

\par\begin{minipage}{\linewidth}\raggedright
% mdparagraph:P265:begin
\textbf{R5.}
% mdparagraph:P265:end

% item: NEW-U03-R05
% mdparagraph:P266:begin
Find both occurrences of geworhte and give their person, number and tense. Identify their objects. Explain why the supplied Bēon and scīnen are not forms of the same indicative pattern.
% mdparagraph:P266:end

\end{minipage}\par\addvspace{4pt}

\par\begin{minipage}{\linewidth}\raggedright
% mdparagraph:P267:begin
\textbf{R6.}
% mdparagraph:P267:end

% item: NEW-U03-R06
% mdparagraph:P268:begin
Analyse gōd in ðæt hit gōd wæs and fēorða in se fēorða dæg. Explain their agreement and their different declensions. Use the supplied closing-formula note for the larger clause.
% mdparagraph:P268:end

\end{minipage}\par\addvspace{4pt}

\reserve{7\baselineskip}
% mdparagraph:P269:begin
\section*{15. Review}
% mdparagraph:P269:end

% mdparagraph:P270:begin
Use the taught forms and ordinary constructions independently. In X2, asterisks mark proposed errors to correct, not reconstructed historical forms. X3 permits more than one defensible composition; keep its stated meaning and referents clear.
% mdparagraph:P270:end

\reserve{7\baselineskip}
% mdparagraph:P271:begin
\subsection*{Review — X1–X3}
% mdparagraph:P271:end

\par\begin{minipage}{\linewidth}\raggedright
% mdparagraph:P272:begin
\textbf{X1.}
% mdparagraph:P272:end

% item: NEW-U03-X01
% mdparagraph:P273:begin
Translate “This good man gives my birds seed. My birds live in his small dwelling. The dwelling is greater than that tree.”
% mdparagraph:P273:end

\end{minipage}\par\addvspace{4pt}

\par\begin{minipage}{\linewidth}\raggedright
% mdparagraph:P274:begin
\textbf{X2.}
% mdparagraph:P274:end

% item: NEW-U03-X02
% mdparagraph:P275:begin
Correct the four marked adjective or numeral forms so that the sentences mean “The second day was good. My birds live in the great dwelling. The bird is smaller.”: Se *ōðra dæg wæs gōd. *Mīn fuglas libbað on ðǣm *micel hofe. Se fugol is *lǣss.
% mdparagraph:P275:end

\end{minipage}\par\addvspace{4pt}

\par\begin{minipage}{\linewidth}\raggedright
% mdparagraph:P276:begin
\textbf{X3.}
% mdparagraph:P276:end

% item: NEW-U03-X03
% mdparagraph:P277:begin
Write four connected Old English sentences: introduce two dwellings of unequal size, compare them, say that an old man has three birds, and give a direct quotation in which he says that he gives his birds food. Use an agreeing possessive in the quotation.
% mdparagraph:P277:end

\end{minipage}\par\addvspace{4pt}

\clearpage
% mdparagraph:P278:begin
\section*{16. Source and editorial notes}
% mdparagraph:P278:end

\begingroup\fontsize{10}{12.8}\selectfont
% mdparagraph:P279:begin
The grammar follows Cook, \textbf{A First Book in Old English}, second edition (1895), §§55–70, 77–80 and 82–85, with the ordinary adjective/\allowbreak{}adverb uses in §§179–182. Wright’s grammar corroborates stem variation, possessives and comparison. The bounded nominative comparison is supported by Fulk, §35. The tables distinguish selected teaching forms from supported alternatives. Forms expanded from a checked stem and paradigm need not be separately attested in a printed source.
% mdparagraph:P279:end

% mdparagraph:P280:begin
Cook’s Exercise I.4 is on printed p.2, one-based PDF p.10 of the selected exercise scan. Its original item wording and references are retained. The source reading is \textbf{First Book}, Reader I, printed p.125, editorially counted lines 8–15, and p.126, lines 1–4; one-based PDF pp.141–142. On p.125 the marginal line numbers 10 and 15 are printed; the other line numbers here are editorial counts. The new reading spans the whole fourth day.
% mdparagraph:P280:end

% mdparagraph:P281:begin
The reading removes apparatus cues, joins line and page breaks, and rejoins the line-divided words tācnum and līehtinge. It otherwise preserves the source’s letters, vowel marks and punctuation, including the colon after twā miclu lēoht. The following eighteen entries preserve Cook’s apparatus for this span. Questions in these historical notes are not additional assigned exercises; use the local support above for forms beyond this unit.
% mdparagraph:P281:end

% mdtable:T29:begin
\par\begingroup
\fontsize{10.5}{13.2}\selectfont
\setlength{\LTleft}{0pt}\setlength{\LTright}{\fill}
\begin{longtable}{@{}>{\raggedright\arraybackslash}p{31.07mm}>{\raggedright\arraybackslash}p{90.12mm}@{}}
\toprule
\textbf{Printed page/\allowbreak{}note} & \textbf{Cook’s note} \\
\midrule
\endfirsthead
\toprule
\textbf{Printed page/\allowbreak{}note} & \textbf{Cook’s note} \\
\midrule
\endhead
\midrule
\multicolumn{2}{r}{\footnotesize Continued on the following page.}\\
\endfoot
\bottomrule
\endlastfoot
125/\allowbreak{}15 & Lat. autem. \\
125/\allowbreak{}16 & See 166. \\
125/\allowbreak{}17 & Gen. sing. \\
125/\allowbreak{}18 & See tōdǣlan. \\
125/\allowbreak{}19 & See tācen, and 24. \\
125/\allowbreak{}20 & See tīd, and 24. \\
125/\allowbreak{}21 & See dæg, and 24. \\
125/\allowbreak{}22 & See gēar. \\
125/\allowbreak{}23 & See 193. a. Write the opt. pret. plur. of this verb. \\
125/\allowbreak{}24 & See twēgen. \\
125/\allowbreak{}25 & See micel. \\
125/\allowbreak{}26 & See 66. \\
125/\allowbreak{}27 & What is the relation of the stem-vowel of this word to that of lēoht? \\
125/\allowbreak{}28 & For niht, instead of nieht, see 19. See 153. d. \\
125/\allowbreak{}29 & See gesęttan, and 189, note. \\
126/\allowbreak{}1 & Opt. pret. = Lat. lucerent. What would be the opt. pres.? \\
126/\allowbreak{}2 & See 156. f. \\
126/\allowbreak{}3 & See 78. \\
\end{longtable}
\par\endgroup

% mdtable:T29:end
% mdparagraph:P282:begin
In these notes “opt.” means optative, Cook’s term for the forms commonly called subjunctive. The supplied present/\allowbreak{}past contrast is sufficient for this reading. His historical vowel question does not require an additional sound-change exercise here.
% mdparagraph:P282:end

\endgroup
\end{document}
